\documentclass[reqno, 11pt]{article}

\usepackage{enumerate}
\usepackage{needspace}
\usepackage{latexsym}
\usepackage{amsmath}
\usepackage{amsfonts}
\usepackage{amsthm}
\usepackage{amssymb,mathtools}
\setcounter{MaxMatrixCols}{20}
\usepackage{newlfont}
\usepackage{graphics}
\usepackage{color}
\usepackage{float}
\usepackage{diagbox}
\usepackage{tocloft}
\usepackage{titlesec}
\usepackage{booktabs,array,longtable,enumitem,microtype}
\usepackage{extpfeil}
\usepackage{centernot}

\usepackage[linesnumbered,ruled,vlined]{algorithm2e}
\usepackage{url}

\usepackage[margin=1in]{geometry}

\usepackage{tabularx}
\usepackage[T1]{fontenc}
\usepackage{lmodern}
\usepackage{rotating}
\usepackage{multirow}
\usepackage{extarrows}
\usepackage[sort,compress,numbers]{natbib}
\usepackage[utf8]{inputenc}
\usepackage{xcolor}
\usepackage{listings}
\usepackage{aliascnt}
\usepackage[pagebackref,colorlinks=true,linkcolor=blue,citecolor=red,urlcolor=blue]{hyperref}
\usepackage[nameinlink,noabbrev,capitalize]{cleveref}
\numberwithin{equation}{section}

\setlist[enumerate]{label=\textup{(\roman*)},leftmargin=*}
\newtheorem{theorem}{Theorem}[section]
\newaliascnt{proposition}{theorem}
\newtheorem{proposition}[proposition]{Proposition}
\aliascntresetthe{proposition}
\newaliascnt{lemma}{theorem}
\newtheorem{lemma}[lemma]{Lemma}
\aliascntresetthe{lemma}
\newaliascnt{corollary}{theorem}
\newtheorem{corollary}[corollary]{Corollary}
\aliascntresetthe{corollary}
\newtheorem{maintheorem}{Theorem}
\renewcommand{\themaintheorem}{\Alph{maintheorem}}
\theoremstyle{definition}
\newaliascnt{definition}{theorem}
\newtheorem{definition}[definition]{Definition}
\aliascntresetthe{definition}
\newaliascnt{remark}{theorem}
\newtheorem{remark}[remark]{Remark}
\aliascntresetthe{remark}
\newaliascnt{example}{theorem}
\newtheorem{example}[example]{Example}
\aliascntresetthe{example}
\newaliascnt{conjecture}{theorem}
\newtheorem{conjecture}[conjecture]{Conjecture}
\aliascntresetthe{conjecture}
\newaliascnt{question}{theorem}
\newtheorem{question}[question]{Question}
\aliascntresetthe{question}
\newaliascnt{problem}{theorem}
\newtheorem{problem}[problem]{Problem}
\aliascntresetthe{problem}
\newaliascnt{alg}{theorem}
\newtheorem{alg}[alg]{Algorithm}
\aliascntresetthe{alg}
\newtheorem{fig}{Figure}

\allowdisplaybreaks[4]

\SetKwInput{KwInput}{Input}                
\SetKwInput{KwOutput}{Output}              

\newcommand{\R}{\mathbb R}
\newcommand{\Z}{\mathbb Z}
\DeclareMathOperator{\vol}{vol}
\DeclareMathOperator{\conv}{conv}
\DeclareMathOperator{\Vol}{Vol}
\DeclareMathOperator{\lw}{lw}
\DeclareMathOperator{\intr}{int}

\setcounter{tocdepth}{1}
\renewcommand{\cftsecfont}{\normalfont\bfseries\color{blue}} 
\renewcommand{\cftsecpagefont}{\normalfont\bfseries\color{blue}} 
\renewcommand{\cftsecleader}{\cftdotfill{\cftdotsep}} 


\title{Proof of the Zaks--Perles--Wills volume conjecture}

\author{Feihu Liu$^{\color{blue} \dag}$, Jinlong Tang$^{\color{blue} \ddag}$, Sihao Tao$^{\color{blue} \P}$, and Zihao Zhang$^{\color{blue} \S}$
\\[2mm]
{\small $^{\color{blue} \dag}$ Center for Combinatorics, LPMC,}\\[-0.8ex]
{\small Nankai University, Tianjin 300071, P.R.~China}\\
{\small $^{\color{blue} \ddag, \P}$ School of Mathematical Sciences,}\\[-0.8ex]
{\small Capital Normal University, Beijing, 100048, P.R.~China}\\
{\small $^{\color{blue} \S}$ School of Mathematics and Statistics,}\\[-0.8ex]
{\small Beijing Institute of Technology, Beijing 102400, P.R.~China}\\
{\small {\color{blue} $^\dag$} Email address: \url{liufeihu7476@163.com}}\\
{\small {\color{blue} $^\ddag$} Email address: \url{jinlong\_tang@cnu.edu.cn}}\\
{\small {\color{blue} $^\P$} Email address: \url{sihao\_tao@cnu.edu.cn}}\\
{\small {\color{blue} $^\S$} Email address: \url{zihao-zhang@foxmail.com}}\\
}


\date{\today}


\begin{document}

\maketitle

\begin{abstract}
The Zaks--Perles--Wills volume conjecture concerns the largest volume
of a lattice polytope with a prescribed number of interior lattice
points. We prove that for $d\ge4$ and $k\ge1$, a $d$-dimensional
lattice polytope with $k$ interior lattice points has normalized volume
at most $(k+1)(s_d-1)^2$, where $s_d$ is the $d$th Sylvester number.
Equality holds if and only if the polytope is affinely unimodularly
equivalent to the Sylvester simplex. Together with the known results
in dimension three, this settles the conjecture for all $d\ge3$.
We prove a volume bound for rational polytopes by induction on the
dimension. At equality, an edge of a lower-dimensional fiber lifts to
an edge of the polytope. Its lattice length identifies the maximizing
polytope.
\end{abstract}

\noindent
\begin{small}
\emph{2020 Mathematics Subject Classification}: Primary 52B20; Secondary 52A38, 52C07.
\end{small}

\noindent
\begin{small}
\emph{Keywords}: Lattice polytope, interior lattice point, volume, Sylvester sequence, lattice projection.
\end{small}

%\tableofcontents

\section{Introduction}

Let $N$ be a lattice of rank $d$, that is, the group of integer
linear combinations of a basis of a $d$-dimensional real vector
space. We write $N_{\R}=N\otimes_{\Z}\R$ for this space.
A \emph{lattice polytope} is the convex hull of finitely many
points of $N$. A $d$-dimensional \emph{simplex} is the convex hull of
$d+1$ affinely independent points. For a full-dimensional polytope
$P\subset N_{\R}$, let $\intr P$ denote its interior and $|P|$ its
volume, measured so that a fundamental parallelepiped of $N$ has
volume one. We use
\[
 I(P):=\#(\intr P\cap N),\qquad \Vol_d(P):=d!\,|P|.
\]
Thus $\Vol_d(P)$ is the normalized volume. In lattice coordinates,
$\Vol_d(P)$ is $d!$ times the usual Euclidean volume. Two lattice polytopes
$P$ and $Q$ are \emph{affinely unimodularly equivalent}, denoted by
$P\simeq Q$, if $Q=UP+b$ in lattice coordinates, where
$U$ is an integer matrix of determinant $1$ or $-1$ and $b$ is a lattice vector.
Such a map preserves both $I(P)$ and $\Vol_d(P)$.

Hensley~\cite[Theorem~3.6]{Hensley} proved that the volume of a
lattice polytope is bounded when its dimension and its positive
number of interior lattice points are fixed. Lagarias and
Ziegler~\cite[Theorem~1]{LZ} obtained an explicit bound for
lattice polytopes with a prescribed positive number of interior
points in a sublattice of the form $u\Z^d$, for a fixed positive
integer $u$. Which lattice polytopes attain the largest possible
volume? The simplices constructed by Zaks, Perles, and Wills~\cite{ZPW}
provide candidates for each dimension and each positive number of
interior lattice points.

The \emph{Sylvester sequence} is defined by $s_1=2$ and
$s_{j+1}=1+\prod_{i=1}^{j}s_i$ for $j\ge1$. Thus its first terms are
$2,3,7,43$. For integers $d\ge3$ and $k\ge1$ and a lattice basis
$e_1,\ldots,e_d$, the \emph{Sylvester simplex} is
\[
 S_{d,k}:=\conv\{0,s_1e_1,\ldots,s_{d-1}e_{d-1},
                           (k+1)(s_d-1)e_d\}.
\]
Zaks, Perles, and Wills~\cite[p.~45]{ZPW} showed that the interior
lattice points in these coordinates are precisely $(1,\ldots,1,j)$
for $1\le j\le k$. The edge vectors from the origin
are the columns of a diagonal matrix. Its determinant is the
normalized volume, so
\[
 \Vol_d(S_{d,k})
  =(k+1)(s_d-1)\prod_{i=1}^{d-1}s_i=(k+1)(s_d-1)^2 .
\]

\Needspace{7\baselineskip}
We use the following formulation of the Zaks--Perles--Wills volume
conjecture due to Balletti and Kasprzyk~\cite[Conjecture~1.5]{BK}.

\begin{conjecture}\label{conj:zpw}
Let $d\ge3$ and $k\ge1$ be integers. Every $d$-dimensional lattice
polytope $P$ with $I(P)=k$ satisfies
$\Vol_d(P)\le\Vol_d(S_{d,k})$. Except when $(d,k)=(3,1)$,
equality holds if and only if $P\simeq S_{d,k}$.
\end{conjecture}

\Cref{tab:known-cases} summarizes previously established cases of
\cref{conj:zpw}. In the table, a lattice
polytope containing the origin in its interior is \emph{reflexive}
if every vertex of its polar
\[
  \{h\in N_{\R}^{*}:h(x)\le1\text{ for all }x\in P\}
\]
belongs to the dual lattice $\operatorname{Hom}(N,\Z)$. Here
$N_{\R}^{*}$ is the space of real linear forms on $N_{\R}$. The two maximizers for
$(d,k)=(3,1)$ are $S_{3,1}$ and $\conv\{0,2e_1,6e_2,6e_3\}$,
up to affine unimodular equivalence. They occur in Kasprzyk's
classification~\cite{Kasprzyk} and are identified explicitly in
the discussion following \cite[Conjecture~1.5]{BK}.

\begin{table}[htbp]
\centering\small
\begin{tabularx}{\textwidth}{@{}ll>{\raggedright\arraybackslash}X@{}}
\toprule
Parameters & Class of polytopes & Reference\\
\midrule
$d=3$, $k=1$ & Lattice polytopes & Kasprzyk~\cite{Kasprzyk}\\
$d=3$, $k=2$ & Lattice polytopes &
  Balletti and Kasprzyk~\cite[Corollary~5.2]{BK}\\
$d=3$, $k\ge3$ & Lattice polytopes &
  Liu, Tang, Tao, and Zhang~\cite[Theorem~1.2]{LTTZ}\\
$d\ge4$, $k=1$ & Lattice simplices &
  Averkov, Kr\"umpelmann, and Nill~\cite[Theorem~2.2]{AKN}\\
$d\ge4$, $k=1$ & Reflexive polytopes &
  Balletti, Kasprzyk, and Nill~\cite[Corollary~1.2]{BKN}\\
$d\ge4$, $k\ge1$ & Lattice simplices &
  Liu, Tang, Tao, and Zhang~\cite[Theorem~1.2]{LTTZII}\\
\bottomrule
\end{tabularx}
\caption{Previously established cases of \cref{conj:zpw}.}\label{tab:known-cases}
\end{table}

Our main result establishes the sharp bound for arbitrary lattice
polytopes of dimension at least four.

\begin{theorem}\label{thm:general-zpw}
Let $d\ge4$ and $k\ge1$ be integers. Every $d$-dimensional lattice
polytope $P$ with exactly $k$ interior lattice points satisfies
\[
                        \Vol_d(P)\le(k+1)(s_d-1)^2 .
\]
Equality holds if and only if $P$ is affinely unimodularly equivalent
to $S_{d,k}$.
\end{theorem}

It follows from \cref{thm:general-zpw} and the three-dimensional
results in \cref{tab:known-cases} that \cref{conj:zpw} holds for
all lattice polytopes of dimension at least three.

Our proof begins with a section through an interior lattice point,
which we translate to the origin. The \emph{lattice width} $\lw(P)$
is the minimum of $\max_P h-\min_P h$ over nonzero linear forms $h$
that take integer values on $N$. Since the origin lies in $\intr P$,
every such form attains a negative minimum and a positive maximum.
Both are integers, since they are attained at vertices of $P$.
Thus $\lw(P)\ge2$. If $\lw(P)=2$, a form attaining the width has
range $[-1,1]$, so every interior lattice point has height zero.

A face of $P$ is the set of maximizers of a linear form on $P$.
Faces of dimensions zero, one, and $d-1$ are called vertices, edges,
and facets, respectively. The relative interior of a face is its
interior in its affine span. Now consider a vertex of the section
at height zero and the smallest face of $P$ containing this vertex.
The vertex lies in the relative interior of that face. If the face had
dimension at least two, its affine span would contain a line of
constant height through the vertex. A short segment of this line on
both sides of the vertex would lie in the section, contrary to
extremality. Thus the vertex is either a vertex of $P$ or lies in an
edge that crosses the section. The endpoints of a crossing edge have
integer heights of opposite sign, hence heights $-1$ and $1$.
Its point at height zero is its midpoint. Consequently, each vertex
$v$ of the section satisfies $2v\in N$. An induction on dimension
must therefore allow rational vertices. For a rational vertex $v$, its
\emph{denominator} is the least positive integer $a$ such that
$av\in N$. We prove a volume bound for rational polytopes in terms
of an upper bound $u$ on these denominators.

For larger widths, suppose the volume of the polytope is at least
the proposed bound. We find a surjection of lattices under which
the image of the polytope has small volume. The volume of this image
bounds the denominators of the vertices
in the fiber over zero. Here the fiber consists of those points of the
polytope that map to zero. At equality, a face of the original polytope
is a translate of a dilation of this fiber. By induction there is
an edge of lattice length $(k+1)(s_d-1)$ in the original lattice
polytope. The lattice-diameter theorem
of Averkov, Kr\"umpelmann, and Nill~\cite[Theorem~2.8]{AKN} determines
the maximizing lattice polytope.

The paper is organized as follows. \Cref{sec:reduction} states the
rational volume bound and relates the volume of a polytope to that
of its fiber over zero. \Cref{sec:tools} establishes the estimates for chord
lengths, covering radii, and planar areas that we use to find a
projection of small volume. \Cref{sec:rational-volumes} proves the
projection bound for $u=2$ and $u\ge6$. For $u=1$, we treat dimension
four in \cref{sec:gzpw-four-chords} and higher dimensions in
\cref{sec:gzpw-wide}. Finally,
\cref{sec:volume-proof} completes the induction and determines the
equality case.

\section{Rational polytopes and their central fibers}\label{sec:reduction}

The induction requires two estimates for the section through an
interior lattice point. One controls its vertex denominators, and
the other relates its volume to the volume of the original polytope.

\subsection{The volume bound for rational polytopes}

A polytope is \emph{rational} if all its vertices have rational
coordinates in a lattice basis. Different vertices may have different
denominators, each bounded by the same integer $u$.
For each positive integer $u$, define $f_0(u)=u$ and
$f_{j+1}(u)=f_j(u)(f_j(u)+1)$ for $j\ge0$. Set
\[
 A_d(u)=\frac{f_{d-1}(u)^2}{d!\,u^{d+1}}\qquad(d\ge1).
\]
The identity $s_{j+2}-1=s_{j+1}(s_{j+1}-1)$ and the initial value
$s_1-1=1$ show that $f_j(1)=s_{j+1}-1$. The value $f_j(u)$ is
obtained by iterating the map $t\mapsto t(t+1)$ exactly $j$ times,
so $f_i(f_j(u))=f_{i+j}(u)$. If the displacement of a segment is
$te$ for a primitive lattice vector
$e$ (a generator of the lattice on its
line), its \emph{lattice length} is $|t|$.

\begin{theorem}\label{thm:general-rational}
Let $C$ be a rational $d$-polytope with $I(C)=k\ge1$, and let $u$ be
a positive integer. Assume that each vertex of $C$ has denominator at
most $u$. If $u=1$ and $d\ge4$, or if $u\in\{2\}\cup\Z_{\ge6}$ and
$d\ge1$, then
\[
                         |C|\le(k+1)A_d(u).
\]
For every such $d,u,k$, there is a polytope attaining equality.
If $(d,u,k)\ne(2,2,1)$ and equality holds, then $C$ is a simplex
with an edge of lattice length $\frac{(k+1)f_{d-1}(u)}{u}$.
\end{theorem}

In dimension one the interval $[0,k+1]$ attains the bound.
For $d\ge2$, take the coordinate simplex of Lagarias and
Ziegler~\cite[Proposition~2.6]{LZ} and divide it by $u$. Their
sublattice $u\Z^d$ becomes $\Z^d$. The resulting simplex has
intercepts $\frac{f_{i-1}(u)+1}{u}$ for $1\le i<d$ and
$\frac{(k+1)f_{d-1}(u)}{u}$ in the last coordinate. Each vertex has
denominator at most $u$. The interior lattice
points are $(1,\ldots,1,j)$ for $1\le j\le k$, by the cited
proposition. By the recurrence,
$\prod_{i=1}^{d-1}(f_{i-1}(u)+1)=\frac{f_{d-1}(u)}{u}$, so the product
of the intercepts is $d!(k+1)A_d(u)$.

\subsection{A volume estimate for the central fiber}

A \emph{convex body} is a compact convex set with nonempty interior
in its ambient space. For a body $K$ in an $m$-dimensional lattice
space, we use the same notation $\Vol_m(K)=m!|K|$ and write $I(K)$
for its number of interior lattice points. We call $K$ \emph{hollow}
if $I(K)=0$.

A \emph{lattice projection} is a linear map $\pi$ induced by a
surjection of lattices. For $y\in\pi(K)$, write $K_y=K\cap\pi^{-1}(y)$.
We use the quotient lattice on the image. On each affine fiber we
use the translate of the volume measure on $\ker\pi$ determined by
the kernel lattice. Above a lattice point of the image, the ambient
lattice induces an affine translate of the kernel lattice on the
fiber space. We call $K_0$ the \emph{central fiber} when $0\in\intr\pi(K)$.
Since a basis of the kernel together with lifts of a
quotient basis forms a basis of the ambient lattice, Fubini's theorem takes the form
$|K|=\int_{\pi(K)}|K_y|\,dy$.

The \emph{relative interior} of a convex set is its interior in
its affine hull. Two elementary facts about relative interiors will be
used below. A linear map takes the relative interior of a convex set
onto the relative interior of its image. Moreover, if an affine
subspace meets the relative interior of a convex set, the relative
interior of their intersection is the intersection of that subspace
with the relative interior of the set
\cite[Corollary~6.5.1 and Theorem~6.6]{Rockafellar}. Since $K$ and
$\pi(K)$ have interior in their respective spaces, the first fact
implies that $\pi^{-1}(y)$ meets $\intr K$ whenever $y\in\intr\pi(K)$.
The second fact then implies
\begin{equation}\label{eq:relative-fiber}
                 \operatorname{relint} K_y=\intr K\cap\pi^{-1}(y).
\end{equation}
In particular, such a fiber has dimension $\dim\ker\pi$.

For nonempty compact convex sets $A,B$ in an $r$-dimensional
space with $r\ge1$, the Brunn--Minkowski inequality~\cite[Theorem~4.1, inequality~(10)]{Gardner}
states that
\[
  |(1-t)A+tB|^{\frac{1}{r}}
    \ge (1-t)|A|^{\frac{1}{r}}+t|B|^{\frac{1}{r}}
    \qquad (0\le t\le1).
\]
Here $A+B$ denotes the
set of all $a+b$ with $a\in A$ and $b\in B$. After identifying fibers
by translation, convexity of $K$ yields
$(1-t)K_y+tK_z\subseteq K_{(1-t)y+tz}$. With $r=\dim\ker\pi\ge1$,
the stated inequality therefore implies that
$y\mapsto |K_y|^{\frac{1}{r}}$ is concave on $\pi(K)$.
It is also upper semicontinuous there. Indeed, a linear splitting of the projection identifies the fibers
with subsets of a common space.
If $y_j\to y$, compactness ensures that every limit point of a
sequence in the fibers $K_{y_j}$ lies in $K_y$. Thus every open
neighborhood of $K_y$ contains $K_{y_j}$ for all sufficiently large
$j$. The volumes of shrinking closed neighborhoods of $K_y$ decrease
to $|K_y|$, so $\limsup_{j\to\infty}|K_{y_j}|\le |K_y|$.

The minimum defining the lattice width is attained. Indeed,
in lattice coordinates a convex body contains a Euclidean ball of
some radius $\varepsilon>0$. An integral linear form $h$ has range
of length at least $2\varepsilon\|h\|$ on this ball. Only finitely
many integral forms have norm below any prescribed bound. A
minimizing nonzero form is primitive, since division by any common
factor of its coefficients would decrease the width. Each nonzero
integral form on the quotient pulls back to a nonzero integral form
on the ambient lattice. Consequently $\lw(\pi(K))\ge\lw(K)$.

For a polytope of width less than $1+\frac{2}{u}$, we use the section through
an interior lattice point in a direction of minimum width. For width
at least $1+\frac{2}{u}$ and $d\ge3$, suppose $|C|\ge(k+1)A_d(u)$. We use
a projection of rank $n$ whose
image has volume at most
\[
                     V_n(u):=\frac{f_n(u)}{n!\,u^{n+1}} .
\]
The following proposition establishes its existence.

\begin{proposition}\label{prop:projections}
Let $C,d,u,k$ be as in \cref{thm:general-rational}.
Assume that $d\ge3$, $\lw(C)\ge1+\frac{2}{u}$, and $|C|\ge(k+1)A_d(u)$.
Then there exists a lattice projection $\pi$ of rank $2\le n<d$
such that $|\pi(C)|\le V_n(u)$.
\end{proposition}

We prove \cref{prop:projections} in
\cref{sec:rational-volumes,sec:gzpw-four-chords,sec:gzpw-wide}.
We first bound the volume in terms of one fiber. Its equality conditions
will determine how an edge of the fiber lifts to the original polytope.

\begin{lemma}\label{lem:central-fiber-volume}
Let $C$ be a $d$-dimensional convex body, and suppose that its image
$B$ under a lattice projection of rank $1\le n<d$ is a polytope with
$0\in\intr B$. Set $r=d-n$, $W=\Vol_n(B)$, and
$\delta=\min_F\Vol_n(\conv(0,F))$, where $F$ ranges over the facets
of $B$. Then
\begin{equation}\label{eq:facet-volume}
                   \Vol_d(C)\le\Vol_r(C_0)\,W(\frac{W}{\delta})^r.
\end{equation}
For each facet $F$, let $\ell_F$ be its affine height, zero on $F$
and one at the origin. At equality, $|C_y|=|C_0|\ell(y)^r$ on $B$,
where $\ell$ is a convex combination of heights $\ell_F$ for which
$\Vol_n(\conv(0,F))=\delta$ and $B$ is a pyramid on $F$, that is,
the convex hull of $F$ and one point outside its affine hull.
If $r>1$, then $\ell=\ell_F$ for one such facet.
\end{lemma}

\begin{proof}
The fiber $C_0$ has positive $r$-dimensional volume by
\eqref{eq:relative-fiber}. Since $0\in\intr B$, the concave function
$y\mapsto |C_y|^{\frac{1}{r}}$ admits an affine upper bound that
agrees with it at zero \cite[Theorem~23.4]{Rockafellar}. After
division by $|C_0|^{\frac{1}{r}}$, this upper bound is an affine
function $\ell$ with $\ell(0)=1$ and
$|C_y|\le |C_0|\ell(y)^r$ on $B$. In particular $\ell\ge0$ on $B$.
The polar of $B$ consists of the linear forms bounded above by one on
$B$. Its vertices are the facet normals normalized to take the value
one on the corresponding facet, namely $1-\ell_F$. Since $1-\ell$
belongs to the polar, it is a convex combination of these normals.
Consequently $\ell$ is a convex combination of the corresponding
facet heights.

For a fixed facet $F$, let $H=\max_B\ell_F$ and
$T=\frac{W}{\Vol_n(\conv(0,F))}$. The pyramid on $F$ with apex at
a point of height $H$ has volume $H|\conv(0,F)|$ and lies in $B$.
Hence $H\le T$. For $0\le t\le H$, write
$v(t)=|\{y\in B:\ell_F(y)\ge t\}|$. Since $\ell_F$ is affine,
the Minkowski convex combination of two such sets lies in the set at
the corresponding intermediate height. By Brunn--Minkowski,
$v(t)^{\frac{1}{n}}$ is concave.

We compute the right derivative at zero. In Euclidean coordinates
chosen so that the lattice has covolume one, the distance between
heights zero and one is $\frac{1}{\|\nabla\ell_F\|}$. The
$(n-1)$-dimensional volumes of the parallel sections at small positive
heights converge to the area of $F$. Indeed, compactness places these
sections in arbitrarily small neighborhoods of $F$ after normal
translation. Conversely, the sections of the inscribed pyramid at
the same heights contain translates of
$\left(1-\frac{t}{H}\right)F$. By integration across the slab of
heights between $0$ and $t$, we obtain
\[
 v(0)=|B|,\qquad v'(0+)=-n|\conv(0,F)|=-\frac{n|B|}{T}.
\]
For $n=1$ this is the corresponding formula for a subinterval.
The right derivative of $v^{\frac{1}{n}}$ at zero is therefore
$-\frac{|B|^{\frac{1}{n}}}{T}$. The graph of a concave function
lies below its right tangent at the left endpoint. Hence
$v(t)\le |B|\left(1-\frac{t}{T}\right)^n$ on $[0,H]$.

The identity $\ell_F(y)^r=\int_0^{\ell_F(y)}rt^{r-1}\,dt$
and Fubini's theorem now yield
\[
 \int_B\ell_F(y)^r\,dy
   =r\int_0^H t^{r-1}v(t)\,dt
   \le r|B|\int_0^T t^{r-1}\left(1-\frac{t}{T}\right)^n\,dt
   =\frac{|B|T^r}{\binom{d}{n}} .
\]
The last equality uses $t=Ts$ and the integer beta integral
$\int_0^1s^{r-1}(1-s)^n\,ds=\frac{(r-1)!\,n!}{(r+n)!}$.
Since $T\le\frac{W}{\delta}$, each facet height has the same integral
upper bound with $T$ replaced by $\frac{W}{\delta}$. The convexity
of the $r$th power on $[0,\infty)$ shows that $\ell^r$ is bounded
pointwise by the same convex combination of the $r$th powers of
the facet heights. Thus
\[
 |C|=\int_B|C_y|\,dy
 \le |C_0|\int_B\ell(y)^r\,dy
 \le \frac{|C_0|\,|B|}{\binom{d}{n}}
               \left(\frac{W}{\delta}\right)^r .
\]
Multiplication by $d!$ proves \eqref{eq:facet-volume} because
$\frac{d!}{\binom{d}{n}}=n!\,r!$.

Suppose equality holds. For every facet height with a positive
coefficient in $\ell$, equality must hold in its integral estimate.
If $H<T$, the integral from $H$ to $T$ of the upper bound has
positive value. If $T<\frac{W}{\delta}$, replacing $T$ by
$\frac{W}{\delta}$ strictly increases the bound. Thus
$H=T=\frac{W}{\delta}$ for each of these facets. Its pyramid with
apex at height $H$ is contained in $B$ and has volume $|B|$.
A proper closed convex subset of a convex body has smaller volume. Indeed, a point of the body outside the subset has a neighborhood disjoint
from the subset which meets the interior of the body. Hence $B$
equals the pyramid. If $r>1$, two distinct facet heights differ on an
open subset of $B$. Strict convexity of the $r$th power would then
make the integral inequality strict, so $\ell$ is a single height.

Finally, equality in the fiber integral forces
$|C_y|=|C_0|\ell(y)^r$ almost everywhere. On $\intr B$ this identity
holds everywhere because a finite concave function is continuous in
the interior of its domain. At a boundary point, an interior sequence
approaching that point and upper semicontinuity imply
$|C_y|\ge |C_0|\ell(y)^r$. The supporting bound yields the reverse
inequality, which completes the proof.
\end{proof}

\subsection{Vertex denominators and equality in the fiber estimate}

We now use the denominators of the vertices to bound both the right
side of \eqref{eq:facet-volume} and the vertex denominators of the central fiber.

\begin{lemma}\label{lem:reduction}
Let $C$ be a rational $d$-polytope with $0\in\intr C$, and assume that
each vertex has denominator at most a positive integer $u$. Let $\pi$
be a lattice projection of rank $1\le n<d$. Assume that one of the
following conditions holds.
\begin{enumerate}
\item $n=1$ and $|\pi(C)|<1+\frac{2}{u}$.
\item $n\ge2$, $|\pi(C)|\le V_n(u)$, and $\lw(\pi(C))>\frac{2}{u}$.
\end{enumerate}
Let $r=d-n$ be the dimension of the central fiber $C_0=C\cap\ker\pi$,
and set $U=f_n(u)$. Each vertex of $C_0$ has denominator at most $U$
in the kernel lattice, and
\begin{equation}\label{eq:one-lift}
            \Vol_d(C)\le \frac{U^{r+1}}{u^{d+1}}\,\Vol_r(C_0).
\end{equation}
If equality holds and $(d,u)\ne(2,1)$, then $\pi(C)$ is a simplex.
For some vertex $v$ of this simplex, $C_v$ is a face of $C$ and a
translate of $\frac{U}{u}C_0$. Moreover, $C$ is the convex hull of $C_v$ and
one point over each of the other vertices of $\pi(C)$.
\end{lemma}

\begin{proof}
Write $B=\pi(C)$ and $W=\Vol_n(B)$. In case~(i), the quotient
lattice has a primitive height coordinate in which $B=[-p,q]$
with $p\le q$, after changing the sign of the coordinate if necessary.
Its endpoints are projected vertices of $C$. Since they are nonzero
rational numbers with denominators at most $u$, both $p$ and $q$ are
at least $\frac{1}{u}$. If $q>1$, the same argument applied to $q-1$
implies $q\ge1+\frac{1}{u}$. This would imply
$p+q\ge1+\frac{2}{u}$, contrary to the hypothesis. We have proved
$\frac{1}{u}\le p\le q\le1$.

Suppose $C$ has $m$ vertices. We write them as
$\frac{z_i}{a_i}$, with $z_i\in N$ and $1\le a_i\le u$ for
$1\le i\le m$. Consider the polytope of possible coefficients
\[
 \left\{c\in\mathbb R_{\ge0}^{m}:
          \sum_{i=1}^{m}c_i(\pi z_i,a_i)=(0,1)\right\}.
\]
The map $c\mapsto\sum_{i=1}^{m}c_i z_i$ takes this polytope onto $C_0$.
Indeed, if $c$ satisfies these equations, the numbers $a_ic_i$ are
nonnegative and sum to one. The convex combination of the vertices
with these coefficients is $\sum_{i=1}^{m}c_i z_i$ and projects to
zero. Conversely, from any convex representation of a point above zero,
division of the $i$th coefficient by $a_i$ produces such a vector $c$.
The coefficient polytope is bounded because $0\le a_ic_i\le1$.

Its vertices map to a finite set whose convex hull is $C_0$, so
each vertex of $C_0$ is the image of at least one of those vertices.
At a vertex $c$, the columns $(\pi z_i,a_i)$ for which $c_i>0$ are
linearly independent. Otherwise a nonzero linear relation supported
on these indices would permit a small perturbation of $c$ in both
directions within the coefficient polytope. The projected vertices
$\frac{\pi z_i}{a_i}$ affinely span the $n$-dimensional image $B$.
Thus the columns $(\frac{\pi z_i}{a_i},1)$ span a space of dimension
$n+1$, as do the original columns. We may therefore extend the
independent columns to $n+1$ independent columns from the same list.
In a lattice basis of the quotient these columns form an integral
square matrix. Cramer's rule, with right side $(0,1)$, shows that
its absolute determinant clears the denominators of $c$, including
all its zero entries. It also clears the denominator of
$\sum_{i=1}^{m} c_i z_i$, since the latter point lies in $\ker\pi$
and its integral multiple belongs to $N\cap\ker\pi$.

For the selected indices $i_0,\ldots,i_n$, the absolute determinant
is
\[
 \left(\prod_{j=0}^{n}a_{i_j}\right)
 \Vol_n\left(\conv\left\{\frac{\pi z_{i_0}}{a_{i_0}},\ldots,
                             \frac{\pi z_{i_n}}{a_{i_n}}\right\}\right).
\]
Indeed, division of column $j$ by $a_{i_j}$ makes its last entry
one. Subtraction of the first column from each remaining column
leaves the usual determinant for the normalized simplex volume.
In case~(ii) the simplex lies in $B$, so the displayed product is at
most $u^{n+1}W\le U$. In case~(i), with selected indices $i,j$,
the absolute determinant is less than $a_i a_j(1+\frac{2}{u})$.
If $a_i=a_j=u$, this determinant is a multiple of $u$ strictly
below $u(u+2)$, and hence at most $u(u+1)$. Otherwise $a_i a_j\le
u(u-1)$ and it is less than $(u-1)(u+2)<u(u+1)$. Since $U=f_1(u)
=u(u+1)$ in this case, the required denominator bound follows.

We next apply \cref{lem:central-fiber-volume} to $C$ and $\pi$. In
case~(i), the two cones from zero to the facets are intervals of
lengths $p,q$. Thus $\delta=p$, $W=p+q$, and
\[
 W\left(\frac{W}{\delta}\right)^r
 =\frac{(p+q)^d}{p^{d-1}}\le\frac{(p+1)^d}{p^{d-1}}
 \le\frac{(u+1)^d}{u}=\frac{U^{r+1}}{u^{d+1}}.
\]
For the second inequality, the logarithmic derivative of
$\frac{(p+1)^d}{p^{d-1}}$ is $\frac{p-(d-1)}{p(p+1)}$.
It is negative on $0<p<1$ for $d\ge2$, and the maximum on
$[\frac{1}{u},1]$ is attained at $\frac{1}{u}$.
In case~(ii), take any facet of $B$ and choose $n$ affinely
independent vertices on it. A linear projection sends vertices to
points of denominator at most $u$, and every vertex of $B$ is the
image of a vertex of $C$. The normalized volume of the simplex
on these $n$ vertices and zero is the absolute determinant of their
integral numerators divided by the product of their denominators.
Its numerator is a nonzero integer because zero lies off the facet
hyperplane. Hence the volume is at least $u^{-n}$. This proves
$\delta\ge u^{-n}$. Together with $W\le\frac{U}{u^{n+1}}$, we obtain
\[
  W\left(\frac{W}{\delta}\right)^r\le
  \frac{U}{u^{n+1}}\left(\frac{U}{u}\right)^r
  =\frac{U^{r+1}}{u^{d+1}} .
\]
In each case, \eqref{eq:facet-volume} now implies \eqref{eq:one-lift}.

Suppose equality holds in \eqref{eq:one-lift}. Then
\eqref{eq:facet-volume} is an equality. Each facet occurring in the
affine upper bound $\ell$ of \cref{lem:central-fiber-volume} is the
base of a pyramid equal to $B$, and its cone from zero has normalized
volume $\delta$. In case~(i), equality requires $p=\frac{1}{u}$ and $q=1$.
In case~(ii), equality requires $\delta=u^{-n}$ and
$W=\frac{U}{u^{n+1}}$. Let $F$ be a facet whose height occurs in
$\ell$ with positive coefficient. Choose $n$ affinely independent
vertices on $F$. Their cone from zero has normalized volume at
least $u^{-n}$ by the determinant bound above. It is contained in
$\conv(0,F)$, whose normalized volume is exactly $\delta=u^{-n}$.
Thus the two cones coincide. Their common intersection with the
hyperplane of $F$ is the simplex on the chosen vertices, so $F$
is a simplex. Since $B$ is a pyramid on $F$, it too is a simplex.
In case~(i), $B$ is already an interval. In both cases the height
of the vertex opposite any facet occurring in $\ell$ with positive
coefficient is $\frac{W}{\delta}=\frac{U}{u}$.

For the equality assertion, assume $(d,u)\ne(2,1)$. We claim that
$\ell=\ell_F$ for one facet $F$. For $r>1$ this follows
from \cref{lem:central-fiber-volume}. Suppose $r=1$. In case~(i) this forces $d=2$, so the exclusion
of $(d,u)=(2,1)$ implies $u>1$.
Hence $p<q$ and only one endpoint attains the minimum $\delta$.
In case~(ii), fix a facet $F$ with positive coefficient in $\ell$.
Since the normalized volume of the cone on $F$ is exactly $u^{-n}$,
the product of the $n$ vertex denominators must be $u^n$
and the numerator determinant must have absolute value one. Each
denominator is therefore $u$, and the integral numerators form a
basis of the quotient lattice. Thus the vertices of $F$ are
$\frac{e_1}{u},\ldots,\frac{e_n}{u}$ for a lattice basis $e_1,\ldots,e_n$. All
coordinates of the opposite vertex $v$ in this basis are negative.
Indeed, the barycentric coordinates of $0$ in the simplex $B$ are
positive, so in every coordinate the contribution of $v$ must
cancel a positive multiple of $\frac{1}{u}$.
If a second facet had positive coefficient, its vertices would be
obtained by replacing some $\frac{e_i}{u}$ by $v$. Its cone determinant
would give $\frac{|v_i|}{u^{n-1}}=u^{-n}$. Hence $v_i=-\frac{1}{u}$, whereas the $i$th coordinates of the
vertices of $F$ are zero or $\frac{1}{u}$. The range of the integral
$i$th coordinate on $B$ would have length $\frac{2}{u}$, contrary to $\lw(B)>\frac{2}{u}$.
This proves the claim.

Let $v$ be the vertex opposite $F$, so that $\ell_F(v)=\frac{U}{u}$.
Choose a point of $C$ above each vertex of $F$, and take their convex
hull with $C_v$. The fiber of this hull at $y\in B$ is a translate
of $(\frac{\ell_F(y)}{\ell_F(v)})C_v$, since this ratio is the barycentric
coordinate of $v$. In fact, $y$ has unique barycentric coordinates
in $B$. All coefficients on the chosen points over $F$ are
therefore fixed, and only the point over $v$ can vary. By \cref{lem:central-fiber-volume},
$|C_v|=|C_0|\ell_F(v)^r$. Thus the fiber of the hull has volume
$|C_0|\ell_F(y)^r=|C_y|$. The hull has the full volume of $C$ and hence
equals $C$. At zero its fiber is a translate of $\frac{u}{U}C_v$. Since a linear form exposing $v$ pulls back
to a form exposing $C_v$, the fiber $C_v$ is a face of $C$ and a
translate of $\frac{U}{u}C_0$.
\end{proof}

\section{Chord lengths and covering radii}\label{sec:tools}

We use the following estimates to find a projection of small volume.
The first two concern the lengths of parallel chords, and the next
two compare lattice-point counts before and after projection.

\subsection{Chord lengths and integrals of their powers}

The \emph{rational diameter} $L(K)$ of a convex body $K$ is the
maximum lattice length of a chord of $K$ over all primitive lattice
directions. If $K$ contains a chord of lattice length at least
$\varepsilon>0$ parallel to a primitive vector $e$, then $e$ belongs to
$\frac{K-K}{\varepsilon}$. This set contains only finitely many lattice
vectors. The body has interior, so at least one direction has a
chord of positive length. We may take $\varepsilon$ equal to this
length and consider just those finitely many directions. In each
direction compactness ensures that the maximum length is attained.

\begin{lemma}\label{aux:chords}
Let $K$ be a convex body of dimension at least two, and let $D$ be
its image under the projection along a primitive lattice vector $e$.
For $y\in D$, let $h(y)$ be the lattice length of the fiber over $y$.
Set $L=\max_D h$ and $D_t=\{y\in D:h(y)\ge t\}$ for $0\le t\le L$.
Then $D_0=D$, $|K|\le L|D|$, and $D_t$ contains a translate of
$(1-\frac{t}{L})D$. In particular, if $n=\dim D$ and $0\le t<L$, then
\[
  \operatorname{lw}(D_t)\ge(1-\frac{t}{L})\operatorname{lw}(D),\qquad
  |D_t|\ge(1-\frac{t}{L})^n|D|.
\]
For every integer $t$ with $0<t<L$, we have $tI(D_t)\le I(K)$.
If $e$ attains the rational diameter of $K$, so that $L=L(K)$, then
\begin{equation}\label{aux:nested}
                   L(D_t)\le\max\{L,2L-2t\}\qquad(0\le t<L).
\end{equation}
\end{lemma}

\begin{proof}
The fiber length $h$ is upper semicontinuous and concave by the
fiber properties in \cref{sec:reduction}. By Fubini's
theorem, $|K|=\int_D h\le L|D|$. Fix $y_*$ with $h(y_*)=L$.
For each $y\in D$, concavity and $h(y)\ge0$ imply
\[
 h\left(\left(1-\frac{t}{L}\right)y+\frac{t}{L}y_*\right)
 \ge\left(1-\frac{t}{L}\right)h(y)+\frac{t}{L}h(y_*)\ge t .
\]
Consequently
\[
              (1-\frac{t}{L})D+(\frac{t}{L})y_*\subseteq D_t.
\]
This proves the bounds on width and volume. Suppose $0<t<L$.
If $y\in\intr D_t$ and $y\ne y_*$, the segment from $y_*$ through
$y$ extends beyond $y$ in $D_t$. Since $h(y_*)>t$, concavity implies
$h(y)>t$. Thus the interior of the fiber over each interior lattice
point of $D_t$ is an open interval of length greater than $t$ in its
lattice line. For integral $t$ this interval contains at least $t$
interior lattice points of $K$. The fibers over distinct lattice points are disjoint, so
$tI(D_t)\le I(K)$.

To prove \eqref{aux:nested}, choose a longest chord of $D_t$ with
displacement $Jv$, where $v$ is primitive in the quotient lattice.
Both fibers of $K$ over the endpoints of this chord have length at
least $t$. Their difference, taken in the order corresponding to $Jv$,
is an interval in the fiber of $K-K$ above $Jv$. The length of the
difference of two intervals is the sum of their lengths, so this
interval has radius at least $t$.
The fiber of $K-K$ over zero contains $[-Le,Le]$. Hence, for every
$0<s<J$, the fiber over $sv$ contains an interval of radius at least
$L(1-\frac{s}{J})+\frac{ts}{J}$ by convexity. Suppose that
$J>\max\{L,2L-2t\}$. At $s=L$ the radius is at least
\[
 L\left(1-\frac{L-t}{J}\right)>\frac{L}{2},
\]
since $J>2(L-t)$. The radius bound depends continuously on $s$,
so it still exceeds $\frac{s}{2}$ for some $s$ with $L<s<J$. Choose a lattice lift
$\widetilde v$ of $v$. Consecutive points $s(\widetilde v+je)$,
$j\in\Z$, are at lattice distance $s$ on this line. Hence one of them
belongs to $K-K$. Since $\widetilde v+je$ projects to $v$, it is
primitive. We obtain a chord of $K$ of length $s>L$, a contradiction.
\end{proof}

We shall apply the following integral bounds both to chord lengths
and to square roots of areas of two-dimensional fibers.

\begin{lemma}\label{aux:support}
Let $h$ be a nonnegative upper semicontinuous concave function on an
$n$-dimensional convex body $D$. Set $D_t=\{y\in D:h(y)\ge t\}$
and $X=\max_D h$. Fix $0<a<X$ and $A\ge |D_a|$, and write $x=\frac{X}{a}$.
There exists a real number $c$ satisfying
\[
 \max\{0,(\frac{|D|}{A})^{\frac{1}{n}}-1\}\le c\le \frac{1}{x-1}
\]
such that, for every positive integer $p$,
\begin{align}
 \int_D h^p
 &\le a^p A\int_0^x pt^{p-1}(1+c(1-t))^n\,dt ,
                         \label{eq:moment-line}\\
 \int_D h^p
 &\le a^p A\frac{(1+c)^{n+p}}{\binom{n+p}{n}c^p}
                    \qquad(c>0). \label{eq:moment-cone}
\end{align}
In particular,
\begin{equation}\label{aux:support-basic}
 \int_D h\le A\max\left\{X,\frac{X^{n+1}}{(n+1)(X-a)^n}\right\}.
\end{equation}
If $x\ge n+1$ and $0<\eta\le (\frac{|D|}{A})^{\frac{1}{n}}-1$, then
\begin{equation}\label{aux:cone}
                 \int_D h\le aA\frac{(1+\eta)^{n+1}}{(n+1)\eta}.
\end{equation}
\end{lemma}

\begin{proof}
By concavity of $h$, we have
$(1-\lambda)D_s+\lambda D_t\subseteq D_{(1-\lambda)s+\lambda t}$
for $0\le\lambda\le1$. The Brunn--Minkowski inequality implies that
$t\mapsto |D_{at}|^{\frac{1}{n}}$ is concave on $[0,x]$. This function is
nonincreasing. Its value at $1$ is positive, since $D_a$ contains a
translate of $(1-\frac{a}{X})D$ by concavity of $h$. Choose a supporting
line at $1$ and write it as $|D_a|^{\frac{1}{n}}(1+c(1-t))$. At $t=0$ its
value is at least $|D|^{\frac{1}{n}}$, and at $t=x$ it is nonnegative. These
two inequalities imply
\[
 \left(\frac{|D|}{A}\right)^{\frac{1}{n}}-1
 \le\left(\frac{|D|}{|D_a|}\right)^{\frac{1}{n}}-1
 \le c\le\frac{1}{x-1}.
\]
Since the function is nonincreasing and $1$ is an interior point of
$[0,x]$, its supporting line has nonpositive slope, so $c\ge0$.
The identity
\[
             \int_D h^p=a^p\int_0^x pt^{p-1}|D_{at}|\,dt
\]
and the supporting bound
$|D_{at}|\le |D_a|(1+c(1-t))^n\le A(1+c(1-t))^n$
imply \eqref{eq:moment-line}. If $c>0$, we may extend the
integral of the supporting bound to $1+\frac{1}{c}$, where the line vanishes.
Under the substitution $t=\frac{1+c}{c}s$ the supporting line
factor $1+c(1-t)$ becomes $(1+c)(1-s)$. Thus the integral up to
$\frac{1+c}{c}$ equals
\[
 \frac{(1+c)^{n+p}}{c^p}\int_0^1 p s^{p-1}(1-s)^n\,ds
 =\frac{(1+c)^{n+p}}{\binom{n+p}{n}c^p},
\]
which proves \eqref{eq:moment-cone}.

For $p=1$ the integral in \eqref{eq:moment-line} is convex as a
function of $c$ on $[0,\frac{1}{x-1}]$. Indeed, for fixed $t\in[0,x]$
the base $1+c(1-t)$ is nonnegative on this entire interval and
the $n$th power is convex there. The value at $c=0$ is $x$.
At $c=\frac{1}{x-1}$ it is
\[
 \int_0^x\left(\frac{x-t}{x-1}\right)^n\,dt
  =\frac{x^{n+1}}{(n+1)(x-1)^n}.
\]
The maximum is attained at an endpoint. Multiplication by $aA$
and the identity $ax=X$ yield \eqref{aux:support-basic}.
If $x\ge n+1$ and $\eta$ satisfies the hypotheses of
\eqref{aux:cone}, then $0<\eta\le c\le\frac{1}{n}$. The logarithmic
derivative of $\frac{(1+c)^{n+1}}{c}$ is
$\frac{nc-1}{c(1+c)}$, which is nonpositive on this interval.
The estimate \eqref{eq:moment-cone} with $p=1$ therefore implies
\eqref{aux:cone}.
\end{proof}

\subsection{Covering radii of successive projections}

The \emph{covering radius} $\mu(K)$ of a convex body $K$ is the least
positive $\lambda$ for which $N+\lambda K=N_\R$. Equivalently, every
translate of $\lambda K$ meets $N$. Indeed, $z+\lambda K$ meets $N$
exactly when $-z\in N+\lambda K$. All the assertions about translates
are invariant under translation of $K$.

If $s>\mu(K)$ and $x_0\in\intr K$, then
$\mu(K)K+(s-\mu(K))x_0\subset s\intr K$ by convexity. Thus every
translate of $s\intr K$ contains a lattice point. If $0<s<\mu(K)$,
some translate of $sK$ misses $N$. At $s=\mu(K)$ there is a translate
of $\mu(K)K$ whose interior misses $N$. To see this, take a sequence
of positive scales increasing to $\mu(K)$ and choose a translate
missing $N$ at each scale. We may take the translation vectors in
a fixed compact fundamental parallelepiped, and pass to a convergent
subsequence. Any lattice point in the interior of the limiting
translate would belong to the preceding translates for all sufficiently
large indices, a contradiction. In particular, a hollow body has
covering radius at least one.

\begin{lemma}\label{aux:torus}
Let $K$ be an $m$-dimensional convex body, and set $L=L(K)$.
Then $|K|\le L^m$. For every real $T\ge L$ and integer $j$ with
$1\le j\le\mu(K)T$, we have
\[
                              j|K|\le T^m.
\]
If $m\ge2$ and $Q$ is the projection along a direction attaining $L$,
then $\mu(Q)\ge\mu(K)-\frac{1}{L}$. More generally, if every translate of
$\lambda\intr Q$ contains at least $j$ lattice points for some
$\lambda>0$ and positive integer $j$, then every translate of
$(\lambda+\frac{h}{L})\intr K$ contains at least $jh$ lattice points for
each positive integer $h$.
\end{lemma}

\begin{proof}
Fix $0<s<\frac{1}{T}$. If $s(K-K)$ contained a nonzero lattice vector,
we could write it as a positive integer multiple of a primitive
vector and obtain a chord of $K$ of lattice length at least
$\frac{1}{s}>L$. Hence the quotient map is injective on $sK$.
The torus obtained by identifying points of $N_\R$ differing by a
lattice vector carries the translation-invariant measure of total mass
one. The image of $sK$ in this torus has measure $s^m|K|$ by the
injectivity just proved. Denote this image by $A$.
For $T=L$, we obtain $|K|\le L^m$ in the limit $s\uparrow\frac{1}{L}$.
For the second bound, the $j$-fold sum of $A$ is the image of
$jsK$, as $K$ is convex. Since $js<\frac{j}{T}\le\mu(K)$,
the translates of $jsK$ by $N$ do not cover $N_\R$. Its image is a
proper compact subset of the torus. Its complement contains a nonempty
open set, so its measure is less than one.
By Macbeath's sum theorem~\cite{Macbeath}, the measure of the sum
of two nonempty compact subsets of the torus is at least the smaller
of one and the sum of their measures. Thus the $j$-fold sum of $A$
has measure at least $\min\{1,js^m|K|\}$. It follows that
$js^m|K|<1$. We obtain $j|K|\le T^m$ in the limit $s\uparrow\frac{1}{T}$.

For the assertions on projections, let $\pi$ denote the projection
onto $Q$ and assume that every translate of $\lambda\intr Q$ contains
at least $j$ lattice points. Let $e$ be a primitive generator of the
kernel lattice. Since $\pi$ projects along a longest chord and the
assertion concerns all translates, we may translate $K$ so that
$[0,Le]$ is such a chord.
For each translation vector $z$, choose $j$ lattice points in
$\pi(z)+\lambda\intr Q$. By \eqref{eq:relative-fiber}, each has a
lift $y\in z+\lambda\intr K$. For every positive integer $h$, we have
$[0,he]\subset(\frac{h}{L})K$, so convexity implies
\[
                       y+[0,he]\subset z+(\lambda+\frac{h}{L})\intr K.
\]
This segment lies on a line projecting to a lattice point, so the
line contains a translate of the kernel lattice. The segment has
lattice length $h$ and therefore
contains at least $h$ lattice points. Indeed, in coordinates with lattice
$\mathbb Z$, every closed interval of integral length $h$ contains
at least $h$ integers. The segments for distinct
projected points are disjoint. Thus every translate of
$(\lambda+\frac{h}{L})\intr K$ contains at least $jh$ lattice points.
With $j=h=1$, we obtain $\mu(K)\le\lambda+\frac{1}{L}$ for every
$\lambda>\mu(Q)$. The desired inequality follows.
\end{proof}

\begin{lemma}\label{aux:flag}
Let $K$ be a $d$-dimensional convex body with $I(K)=k$. Then
\begin{equation}\label{eq:gzpw-density}
                      1-\mu(K)\le(\frac{k}{|K|})^{\frac{1}{d}}.
\end{equation}
Set $Q_d=K$. For $2\le i\le d$, let $Q_{i-1}$ be the projection
of $Q_i$ along a direction attaining $L_i=L(Q_i)$, and write
$\mu_i=\mu(Q_i)$ for $1\le i\le d$.
For every $1\le m<d$ and positive integers $h_{m+1},\ldots,h_d$,
\begin{equation}\label{eq:gzpw-integer-count}
     \prod_{i=m+1}^d h_i\ge k+1
        \quad\Longrightarrow\quad
        1-\mu_m\le\sum_{i=m+1}^d \frac{h_i}{L_i} .
\end{equation}
In particular, if $|K|\ge(k+1)\tau$ for some $\tau>0$ and $d\ge2$, then
\begin{align}
  1-\mu_m &\le (d-m)(\frac{|Q_m|}{\tau})^{\frac{1}{d-m}}
                      +\sum_{i=m+1}^d \frac{1}{L_i} ,
                         \label{eq:gzpw-continuous-count}\\
  1-\mu_{d-1}&\le \frac{k+1}{L_d}\le \frac{|Q_{d-1}|}{\tau} .
                         \label{eq:gzpw-one-step}
\end{align}
\end{lemma}

\begin{proof}
If $\mu(K)\ge1$, then \eqref{eq:gzpw-density} is immediate.
Otherwise fix $\mu(K)<\lambda<1$. The average number of lattice points in a
translate of $(1-\lambda)\intr K$ is $(1-\lambda)^d|K|$. Indeed,
we may average the translation over a fundamental parallelepiped of
volume one. For each lattice point the integral of its membership
indicator is the volume of the part of $(1-\lambda)\intr K$ lying in
the corresponding translate of the reflected parallelepiped. These
parallelepipeds tile the space up to their boundaries, so the integrals
sum to the volume of $(1-\lambda)K$. We choose
a translate $z+(1-\lambda)\intr K$ containing at least that many
lattice points, and choose a lattice point in $-z+\lambda\intr K$.
The sums of this point and the points in the first translate are
distinct interior lattice points of $K$, since
$(1-\lambda)\intr K+\lambda\intr K=\intr K$. Hence
$k\ge(1-\lambda)^d|K|$ for each such $\lambda$, which proves
\eqref{eq:gzpw-density}.

For $s>\mu_m$, every translate of $s\intr Q_m$ contains a lattice
point. By \cref{aux:torus}, passage from $Q_{i-1}$ to $Q_i$
increases the scale by $\frac{h_i}{L_i}$ and multiplies the
guaranteed number of lattice points by $h_i$. After the steps $i=m+1,\ldots,d$, every translate of
\[
               \left(s+\sum_{i=m+1}^d \frac{h_i}{L_i}\right)\intr K
\]
contains at least $\prod_{i=m+1}^d h_i$ lattice points. Suppose
the product is at least $k+1$ and
$\mu_m+\sum_{i=m+1}^{d}\frac{h_i}{L_i}<1$.
For $s$ sufficiently close to $\mu_m$, the scale in this display is
less than one. A dilate of $K$ by this factor about an interior point
lies in $\intr K$, a contradiction. This proves
\eqref{eq:gzpw-integer-count}.

By \cref{aux:chords}, we have $|K|\le |Q_m|\prod_{i=m+1}^d L_i$.
Set $a=(\frac{|Q_m|}{\tau})^{\frac{1}{d-m}}$ and take $h_i=\lceil aL_i\rceil$.
Then $\prod_{i=m+1}^d h_i\ge a^{d-m}\prod_{i=m+1}^d L_i
\ge \frac{|K|}{\tau}\ge k+1$ and $\frac{h_i}{L_i}\le a+\frac{1}{L_i}$. Thus \eqref{eq:gzpw-continuous-count} follows from
\eqref{eq:gzpw-integer-count}. For $m=d-1$, we may take
$h_d=k+1$ directly. The first inequality in \eqref{eq:gzpw-one-step}
follows, and the second follows from $|K|\le L_d|Q_{d-1}|$.
\end{proof}

\subsection{Planar area bounds and their application to fibers}

\begin{lemma}\label{aux:planar}
Let $B$ be a planar convex body with $w=\lw(B)>1$ and $k=I(B)$.
For $1<t\le2$, set $\Phi(t)=\frac{t^2}{2(t-1)}$. If $1<w\le2$, then
$|B|\le(k+1)\Phi(w)$. If $w\ge2$, then $|B|\le2$ for $k=0$
and $|B|\le2k+\frac{5}{2}$ for $k\ge1$. In particular, a hollow planar
body of width at least $\rho$, where $1<\rho\le2$, has area at most
$\Phi(\rho)$. For every lower bound $1<\rho\le w$ on the width,
we also have
\begin{equation}\label{aux:planar-radius}
             \min\{\mu(B),\frac{2}{\rho}\}\le \frac{1}{\rho}+\frac{\rho}{2|B|}.
\end{equation}
\end{lemma}

\begin{proof}
For a hollow planar convex body, Averkov and
Wagner~\cite[Theorems~2.1 and~2.2]{AW} state that
$w\le1+\frac{2}{\sqrt3}$. Their area bound for $1<w\le2$ is
$\Phi(w)$. For $2<w\le1+\frac{2}{\sqrt3}$ it is
\[
 \frac{3w^2}{3w+1-\sqrt{1+6w-3w^2}}
 =\frac{3w+1+\sqrt{1+6w-3w^2}}{4}\le 2 .
\]
The equality follows by rationalizing the denominator. For the last
inequality, $7-3w>0$ on the specified interval and
$(7-3w)^2-(1+6w-3w^2)=12(w-2)^2\ge0$.

For positive $k$, Bohnert~\cite[Theorem~1.2]{Bohnert} proves the
bounds $(k+1)\Phi(w)$ for $1<w\le2$ and $2k+\frac{5}{2}$ for
$2<w\le5$. For $w>5$, his Theorem~1.1 implies $|B|\le2(k+1)$.
At $w=2$, the first bound is $2(k+1)\le2k+\frac{5}{2}$.
The derivative
$\Phi'(t)=\frac{t(t-2)}{2(t-1)^2}$ is nonpositive on $(1,2]$.
If a hollow body has width in $[\rho,2]$, its area is thus at most
$\Phi(\rho)$. For a hollow body of width greater than two, the bound
$2=\Phi(2)\le\Phi(\rho)$ applies.

To deduce \eqref{aux:planar-radius}, fix
$0<s<\min\{\mu(B),\frac{2}{\rho}\}$. A translate of $sB$ is hollow,
and its width is at least $s\rho$. If $s\rho>1$, then $1<s\rho<2$,
so the preceding bound implies
\[
 s^2|B|\le\Phi(s\rho)=\frac{s^2\rho^2}{2(s\rho-1)}.
\]
After cancellation and multiplication by $2(s\rho-1)$ this becomes
$2|B|(s\rho-1)\le\rho^2$, or
$s\le\frac{1}{\rho}+\frac{\rho}{2|B|}$. If $s\rho\le1$, this
inequality follows directly from $s\le\frac{1}{\rho}$. The bound
therefore persists as $s$ increases to the indicated minimum.
\end{proof}

For $1<v\le2$, set $\Phi_+(v)=\max\{\frac{9}{4},\Phi(v)\}$. If a planar body
has width at least $v$ and fewer than $h$ interior lattice points,
where $h$ is a positive integer, then its area is at most $h\Phi_+(v)$.
Indeed, for width greater than two and $k\ge1$,
$2k+\frac{5}{2}\le(\frac{9}{4})(k+1)$, and the other cases follow directly from
\cref{aux:planar}. We use this area estimate for the level sets of
chord lengths in dimension three.

\begin{lemma}\label{lem:three-count}
Let $T$ be a three-dimensional convex body with $I(T)<p$ for a
positive integer $p$. Set $J=L(T)$, and suppose the projection
along a direction attaining $J$ has width at least $w$, with $w\le2$.
For every positive integer $m$ such that $t=\frac{J}{m}\ge\frac{5}{2}$ and
$v=w(1-\frac{1}{t})>1$, we have
\[
 |T|\le J\left(1+\left\lfloor\frac{p-1}{m}\right\rfloor\right)\Phi_+(v).
\]
In particular, $\frac{|T|}{p}\le(t+\frac{J}{p})\Phi_+(v)$. If $m$ divides $p$, then
$\frac{|T|}{p}\le t\Phi_+(v)$.
\end{lemma}

\begin{proof}
Let $D_m$ be the set of points in the projection whose fibers have
length at least $m$. By \cref{aux:chords}, it has width at least $v$
and at most $\lfloor\frac{p-1}{m}\rfloor$ interior lattice points.
Here $t\ge\frac{5}{2}$ implies
$3(1-\frac{1}{t})^2\ge3(\frac{3}{5})^2>1$.
In \eqref{aux:support-basic}, we now take $n=2$, $X=J$, $a=m$ and
$A=|D_m|$. The second expression in the maximum is
$\frac{J}{3(1-\frac{1}{t})^2}\le J$, so $|T|\le J|D_m|$.
The planar area bound applied to $D_m$ therefore proves the first
assertion. Since $\lfloor\frac{p-1}{m}\rfloor\le\frac{p}{m}$
and $\frac{J}{m}=t$, division by $p$ yields
$\frac{|T|}{p}\le(t+\frac{J}{p})\Phi_+(v)$.
If $m$ divides $p$, then $\lfloor\frac{p-1}{m}\rfloor=\frac{p}{m}-1$,
and the exact factor is $\frac{J}{m}=t$.
\end{proof}

\section{Small projections of rational polytopes}\label{sec:rational-volumes}

We assume the hypotheses of \cref{prop:projections} with $u=2$ or $u\ge6$.
Suppose, to the contrary, that every proper lattice projection of rank
$i\ge2$ has volume greater than $V_i(u)$. We write $V_i=V_i(u)$ and
$A_i=A_i(u)$ throughout this section. We start with $Q_d=C$ and project along a longest chord in each
dimension until we reach a planar image
$Q_2$, as in \cref{aux:flag}. Write $L_i=L(Q_i)$ and $\mu_i=\mu(Q_i)$.
For $2\le i<d$, the composite projection is a lattice projection, so
$|Q_i|>V_i$. At the top dimension, $|C|\ge(k+1)A_d\ge2A_d>V_d$,
since $\frac{2A_d}{V_d}=\frac{2f_{d-1}(u)}{f_{d-1}(u)+1}>1$.
It follows from \cref{aux:torus} that $L_i^i\ge|Q_i|>V_i$ for every
$2\le i\le d$.

Every planar lattice image $B$ has width at least $w=1+\frac{2}{u}$.
For $u\ge6$, its area satisfies
\[
 |B|>V_2=\frac{(u+1)(u^2+u+1)}{2u^2}>\frac{u+2}{2}.
\]
The last inequality follows by subtracting $u^2(u+2)$ from the
numerator, which leaves $2u+1$. For $u=2$, twice each vertex of $B$
is integral, so $2!\,|2B|=8|B|$ is an integer. Since $|B|>V_2=\frac{21}{8}$,
we obtain $|B|\ge\frac{22}{8}=\frac{11}{4}$. In either case the right
side of \eqref{aux:planar-radius}, with lower width bound $w$, is less
than one. Since $\frac{2}{w}\ge1$, the minimum in that estimate must be
$\mu(B)$. The two area bounds therefore imply
\begin{equation}\label{eq:rat-planar-radius}
 \begin{aligned}
 \mu(B)&<\frac{u}{u+2}+\frac{1}{u}\le1-\frac{1}{2u} &&(u\ge6),\\
 \mu(B)&\le\frac{1}{2}+\frac{4}{11}=\frac{19}{22} &&(u=2).
 \end{aligned}
\end{equation}
Indeed, $\frac{w}{2|B|}<\frac{1}{u}$ in the first case, and
$\frac{u}{u+2}+\frac{1}{u}\le1-\frac{1}{2u}$ is equivalent to $u\ge6$.

\subsection{The projection bound in dimension three}

\begin{proof}[Proof of \cref{prop:projections} in dimension three]
Set $X=\frac{L_3}{k+1}$. From \eqref{eq:gzpw-one-step},
$1-\mu_2\le\frac{1}{X}$. Thus \eqref{eq:rat-planar-radius} implies
$X<2u$ for $u\ge6$ and $X\le\frac{22}{3}$ for $u=2$.
We use \cref{lem:three-count} with $T=C$ and $p=k+1$. For any
$\frac{5}{2}\le a\le L_3$, the integer $m=\lfloor\frac{L_3}{a}\rfloor$
satisfies $m\ge1$ and $ma\le L_3<(m+1)a$. Consequently,
$a\le t=\frac{L_3}{m}<2a$.

Suppose first that $u\ge6$. Since $f_2(u)=u(u+1)(u^2+u+1)$
and $u^2+u+1>u(u+1)$, we have
\[
 A_3=\frac{(u+1)^2(u^2+u+1)^2}{6u^2}>\frac{(u+1)^4}{6}.
\]
Also, $(u+1)(\frac{u+1}{u+2})^3\ge7(\frac{7}{8})^3>3$. Hence
$L_3^3\ge|C|\ge2A_3>(u+2)^3$, and we may take $a=u+2$.
The width parameter in \cref{lem:three-count} satisfies
\[
 w\left(1-\frac{1}{t}\right)\ge
 \left(1+\frac{2}{u}\right)\left(1-\frac{1}{u+2}\right)=1+\frac{1}{u}.
\]
At this lower endpoint, $\Phi(1+\frac{1}{u})=\frac{(u+1)^2}{2u}>
\frac{9}{4}$, so $\Phi_+$ has the same value. By monotonicity of $\Phi_+$
and the bounds $t<2a$, $X<2u$, it follows that
\[
 \frac{|C|}{k+1}\le(2a+2u)\frac{(u+1)^2}{2u}
    =\frac{2(u+1)^3}{u}<\frac{(u+1)^4}{6}<A_3.
\]
Here the strict comparison reduces to $12<u(u+1)$. This contradicts
the volume hypothesis.

Suppose next that $u=2$. Then $w=2$, $A_3=\frac{147}{8}$, and
$L_3^3\ge2A_3>(\frac{5}{2})^3$. For $t>2$ we have
\[
 \Phi\left(2-\frac{2}{t}\right)=2+\frac{2}{t(t-2)},\qquad
 t\Phi\left(2-\frac{2}{t}\right)=2t+\frac{2}{t-2}.
\]
Both functions are convex. For the first one, $\frac{1}{t(t-2)}$
has convex logarithm $-\log t-\log(t-2)$ and is therefore convex.
For the second one this follows directly by differentiation.
Thus $t\Phi_+(2-\frac{2}{t})$ and $(t+\frac{5}{2})\Phi_+(2-\frac{2}{t})$
are convex as well. In each case replacing $\Phi$ by $\Phi_+$ amounts
to taking the maximum with an affine function of $t$.

If $X\le\frac{5}{2}$, then $\frac{5}{2}\le t\le5$ by our choice $a=\frac{5}{2}$. The lemma bounds $\frac{|C|}{k+1}$ by
$(t+\frac{5}{2})\Phi_+(2-\frac{2}{t})$. At the two endpoints
$\Phi_+$ is respectively $\frac{18}{5}$ and $\frac{9}{4}$, so the
corresponding bounds are $18$ and $\frac{135}{8}$. If instead
$\frac{5}{2}\le X\le\frac{22}{3}$, take $m=k+1$. Now $m$ divides $p$
and $t=X$, so the sharper bound of the lemma is $t\Phi_+(2-\frac{2}{t})$.
Its values at $\frac{5}{2}$ and $\frac{22}{3}$ are $9$ and $\frac{33}{2}$.
A convex function on an interval is bounded by its larger endpoint
value. All four values are less than $\frac{147}{8}$, again a contradiction.
\end{proof}

\subsection{Three-dimensional images in the half-integral case}

For $u=2$ and $d\ge4$, the three-dimensional image has volume
$|Q_3|>V_3(2)=\frac{301}{16}$. Each of its planar lattice images
is also an image of $C$, so its area is at least $\frac{11}{4}$.

\begin{lemma}\label{lem:half-three-image}
Let $Q$ be a three-dimensional convex body with $|Q|>\frac{301}{16}$.
Suppose that each planar lattice projection of $Q$ has width at
least two and area at least $\frac{11}{4}$. Then
\begin{equation}\label{eq:rat-half-radius-three}
                           \mu(Q)<\frac{15}{16}.
\end{equation}
If a three-dimensional convex body $T$ in the same lattice space
contains a translate of $\frac{9}{10}Q$ and has fewer than $p$ interior
lattice points for a positive integer $p$, then
\begin{equation}\label{eq:rat-thick-three}
                               |T|\le88p.
\end{equation}
\end{lemma}

\begin{proof}
Suppose $\mu(Q)\ge\frac{15}{16}$. There exists a hollow translate
$H$ of $\frac{15}{16}Q$, so
$|H|>(\frac{15}{16})^3\frac{301}{16}>15$. Write $J=L(H)$ and take
the sets $D_t$ of \cref{aux:chords} in a direction attaining $J$.
If $J\le3$, we can apply \cref{aux:torus} with $T=j=3$, since
$\mu(H)\ge1$. This implies $|H|\le\frac{3^3}{3}=9$.

If $3\le J\le7$, the set $D_1$ is hollow. Its width is at least
$\frac{15}{8}(1-\frac{1}{J})$ by \cref{aux:chords}, since the projected
body has width at least $\frac{15}{16}\cdot2$. This lower bound lies
between $1$ and $2$. Since $3(1-\frac{1}{J})^2\ge1$,
\eqref{aux:support-basic} with $a=1$ implies $|H|\le J|D_1|$. By
\cref{aux:planar},
\[
 |H|\le J\Phi\left(\frac{15}{8}\left(1-\frac{1}{J}\right)\right)
       =\frac{225(J-1)^2}{16(7J-15)}.
\]
On this interval the denominator is positive. The function is convex,
as one sees by writing $J-1=\frac{1}{7}(7J-15)+\frac{8}{7}$ and
expanding the square. The endpoint values are $\frac{75}{8}$ and
$\frac{2025}{136}$, both less than $15$.

If $J\ge7$, the width of $D_1$ is at least
$\frac{15}{8}(1-\frac{1}{7})>\frac{8}{5}$, whence
$|D_1|\le\Phi(\frac{8}{5})=\frac{32}{15}$.
The full projection satisfies
\[
 |D_0|\ge\left(\frac{15}{16}\right)^2\frac{11}{4}>
           \left(\frac{17}{16}\right)^2\frac{32}{15}.
\]
Thus $\sqrt{\frac{15|D_0|}{32}}-1>\frac{1}{16}$, and $J\ge3$.
These are the two hypotheses of \eqref{aux:cone} for $n=2$,
$a=1$, $A=\frac{32}{15}$, and $\eta=\frac{1}{16}$. Hence
\[
 |H|\le\frac{32}{15}\left(\frac{17}{16}\right)^3
              \frac{16}{3}=\frac{4913}{360}<15.
\]
The alternatives exhaust all $J$ and contradict $|H|>15$.
This proves \eqref{eq:rat-half-radius-three}.

For the second assertion, write $J=L(T)$ and use the sets $D_t$ for
$T$ in a direction attaining $J$. Denote the projection of $Q$ along
this direction by $B$. From \eqref{aux:planar-radius} with width bound
two, $\mu(B)\le\frac{1}{2}+\frac{4}{11}=\frac{19}{22}<1$.
If $J>p$, then $pI(D_p)\le I(T)<p$, so $D_p$ is hollow. It contains
a translate of $\frac{9}{10}(1-\frac{p}{J})B$ by \cref{aux:chords}.
This scale cannot exceed $\mu(B)$, since above the covering radius
every translate contains an interior lattice point. Consequently,
\[
 \frac{9}{10}\left(1-\frac{p}{J}\right)\le\frac{19}{22},
 \qquad \frac{p}{J}\ge\frac{4}{99},\qquad J\le\frac{99p}{4}<25p .
\]
If $J\le p$, then $J<25p$ also holds. For $J<4$, the bound
$|T|\le J^3<64\le88p$ suffices. For $J\ge4$ we take
$m=\lfloor\frac{J}{4}\rfloor$ and $t=\frac{J}{m}$, so $4\le t<8$.
The planar image of $T$ has width at least $\frac{9}{10}\cdot2$.
Hence the width parameter in \cref{lem:three-count} is at least
$\frac{9}{5}\cdot\frac{3}{4}=\frac{27}{20}>\frac{4}{3}$ and at most
two if we use $\frac{9}{5}$ as the width bound. Since $\Phi_+$ is
nonincreasing and $\Phi_+(\frac{4}{3})=\frac{8}{3}$, it follows that
\[
 \frac{|T|}{p}\le\left(t+\frac{J}{p}\right)\Phi_+\left(\frac{4}{3}\right)
       \le(8+25)\frac{8}{3}=88.\qedhere
\]
\end{proof}

\subsection{The projection bound in higher dimensions}

We first record the growth estimate used in this section and in
\cref{sec:gzpw-wide}. For every $u\ge1$, both sequences $M_i=V_i(u)$
and $M_i=A_i(u)$ satisfy
\begin{equation}\label{eq:volume-scale-growth}
 \frac{M_{i+1}}{M_i^2}\ge\frac{u^i i!}{i+1}\ge1\qquad(i\ge3).
\end{equation}
Indeed, substitution of their definitions and of the recurrence for
$f_i$ gives the two quotients in the form
\[
 \frac{V_{i+1}(u)}{V_i(u)^2}=\frac{u^i i!}{i+1}
                 \left(1+\frac{1}{f_i(u)}\right),\qquad
 \frac{A_{i+1}(u)}{A_i(u)^2}=\frac{u^i i!}{i+1}
                 \left(1+\frac{1}{f_{i-1}(u)}\right)^2 .
\]
Thus $M_{i+1}\ge M_i^2$. If $\sqrt[i]{M_i}>4$ and $i\ge3$, then
\[
 \frac{\sqrt[i+1]{M_{i+1}}}{\sqrt[i]{M_i}}
     \ge\bigl(\sqrt[i]{M_i}\bigr)^{\frac{i-1}{i+1}}>2,
\]
because $\frac{i-1}{i+1}\ge\frac{1}{2}$. The same hypothesis holds
at $i+1$, so the roots at least double at every subsequent step.

\begin{proof}[Proof of \cref{prop:projections} in dimensions $d\ge4$]
Set $T_i=\sqrt[i]{V_i}$. By \cref{aux:torus},
$\mu_{i-1}\ge\mu_i-\frac{1}{L_i}$, and $L_i>T_i$.
At $i=d$, \eqref{eq:gzpw-density} and $L_d^d\ge|C|\ge(k+1)A_d$ imply
\begin{equation}\label{eq:rat-top-loss}
 1-\mu_{d-1}\le \sqrt[d]{\frac{k}{|C|}}+\frac{1}{L_d}
 \le\frac{\sqrt[d]{k}+1}{\sqrt[d]{(k+1)A_d}}
 \le\sqrt[d]{\frac{2^{d-1}}{A_d}}<\frac{2}{T_d}.
\end{equation}
The penultimate inequality follows from convexity of the $d$th power,
and the last one follows from $V_d<2A_d$. For $3\le j\le d$ with $T_j>4$,
we claim that $1-\mu_{j-1}<\frac{2}{T_j}$. By
\eqref{eq:volume-scale-growth}, $T_{i+1}\ge2T_i$ for every $i\ge j$.
Thus whenever $j\le i<d$ and $1-\mu_i<\frac{2}{T_{i+1}}$, we obtain
\[
 1-\mu_{i-1}\le1-\mu_i+\frac{1}{L_i}
 <\frac{2}{T_{i+1}}+\frac{1}{T_i}\le\frac{2}{T_i}.
\]
The initial estimate \eqref{eq:rat-top-loss} therefore proves the claim by descending induction.

Suppose $u\ge6$. Since $f_1(u)=u(u+1)$ and $f_{i+1}(u)>f_i(u)^2$,
we have $f_i(u)\ge[u(u+1)]^{2^{i-1}}$ for every $i\ge1$. In particular,
\[
 \frac{T_4}{u}\ge\frac{(u+1)^2}{\sqrt[4]{24u}}>14,\qquad
 V_3\ge\frac{(u+1)^4}{6}>7^3 .
\]
The first lower bound equals $\frac{49}{\sqrt{12}}>14$ at $u=6$ and
increases with $u$, since its logarithmic derivative is
$\frac{2}{u+1}-\frac{1}{4u}>0$. The second follows from
$7^4>6\cdot7^3$. The claim with $j=4$ now implies
$1-\mu_3<\frac{1}{7u}$, while $L_3>\sqrt[3]{V_3}>7$. It follows that
\[
 \mu_3-\frac{1}{L_3}>1-\frac{1}{42}-\frac{1}{7}>\frac{3}{4}.
\]
Consequently $\lfloor\mu_3 L_3\rfloor>\frac{3}{4}L_3$. By
\cref{aux:torus} with $T=L_3$ and $j=\lfloor\mu_3 L_3\rfloor$,
\[
 L_3^3\ge\lfloor\mu_3 L_3\rfloor |Q_3|>
             \frac{3}{4}L_3 V_3\ge \frac{L_3(u+1)^4}{8}.
\]
Thus $L_3>\frac{(u+1)^2}{\sqrt{8}}>2\sqrt{2}\,u>\frac{14u}{5}$,
where $(u+1)^2>8u$ for $u\ge6$. These estimates yield
\[
 1-\mu_2\le1-\mu_3+\frac{1}{L_3}<
         \frac{1}{7u}+\frac{5}{14u}=\frac{1}{2u},
\]
contrary to \eqref{eq:rat-planar-radius}.

Suppose now that $u=2$. By the observations before
\cref{lem:half-three-image}, that lemma applies to $Q_3$.
If $d=4$, then $|C|\ge(k+1)A_4$, where $A_4=\frac{1806^2}{768}>4200$.
By \eqref{eq:gzpw-density}, $1-\mu_4\le\sqrt[4]{\frac{k}{|C|}}
<\frac{1}{\sqrt[4]{4200}}<\frac{1}{2}$. If $L_4\le10$, the choices
$T=10$ and $j=2$ in \cref{aux:torus} would yield
$|C|\le\frac{10^4}{2}<2A_4$. Hence $L_4>10$. Also,
$\frac{1}{16}<1-\mu_3\le\frac{k+1}{L_4}$ by
\eqref{eq:rat-half-radius-three} and \eqref{eq:gzpw-one-step}.

Set $n=\lfloor\frac{L_4}{10}\rfloor$. Then $n\ge1$ and
$10\le\frac{L_4}{n}<20$. Let $T$ be the set of points of $Q_3$
whose fibers have length at least $n$. By \cref{aux:chords}, this
set contains a translate of $(1-\frac{n}{L_4})Q_3$, hence also one
of $\frac{9}{10}Q_3$, and $I(T)\le\lfloor\frac{k}{n}\rfloor$.
It follows from \eqref{eq:rat-thick-three} that
$|T|\le88(\lfloor\frac{k}{n}\rfloor+1)$. Since
$4(1-\frac{n}{L_4})^3\ge4(\frac{9}{10})^3>1$,
\eqref{aux:support-basic} implies $|C|\le L_4|T|$. Therefore
\[
 \frac{|C|}{k+1}\le88\left(\frac{k}{k+1}\frac{L_4}{n}+
                              \frac{L_4}{k+1}\right)
        <88(20+16)<4200,
\]
which contradicts $|C|\ge(k+1)A_4$.

If $d\ge5$, then $V_4=\frac{3263442}{768}>4200$ and
$T_5^5=\frac{3263442\cdot3263443}{5!\,2^6}>65^5$ by the recurrence
for $f_i(2)$. The claim with $j=5$ implies
$\mu_4>1-\frac{2}{65}>\frac{31}{32}$. Since
$\mu_4-\mu_3\le\frac{1}{L_4}$ and $\mu_3<\frac{15}{16}$, we obtain
$\frac{1}{L_4}>\frac{31}{32}-\frac{15}{16}=\frac{1}{32}$.
For $L_4\le15$ we may take $T=15$, $j=14$ in \cref{aux:torus},
because $15\mu_4>14$. This yields $|Q_4|\le\frac{15^4}{14}<4200$.
For $15<L_4<32$, we have $\mu_3>\frac{31}{32}-\frac{1}{15}>
\frac{9}{10}$. There is a hollow translate of $\frac{9}{10}Q_3$.
Its volume is at most $88$ by \eqref{eq:rat-thick-three} with $p=1$.
We conclude that
\[
 |Q_4|\le L_4|Q_3|<
       32\cdot88\left(\frac{10}{9}\right)^3<4200.
\]
Both alternatives contradict $|Q_4|>V_4$. This completes the proof.
\end{proof}

\section{Small projections of four-dimensional lattice polytopes}\label{sec:gzpw-four-chords}

For a lattice four-polytope $P$, the thresholds in \cref{prop:projections}
are $V_2(1)=3$ and $V_3(1)=7$. Every lattice projection of $P$ is a
lattice polytope. In dimension $i$, the volume of such a polytope
multiplied by $i!$ is an integer, as one sees by triangulating it
into simplices with lattice vertices and taking determinants.
If all images of ranks two and three exceed their respective thresholds,
their normalized volumes are therefore at least $7$ and $43$.
We shall prove that these lower bounds force
$|P|<\frac{147}{2}(I(P)+1)$ whenever $\lw(P)\ge3$ and $I(P)\ge1$.
The proof uses the lattice points in successive fibers of $P$.

\subsection{The volume of a hollow three-dimensional body}

\begin{lemma}\label{aux:hollow-envelope}
Let $K$ and $Q$ be convex bodies in the same three-dimensional lattice
space. Suppose that $K$ is hollow and contains a translate of $sQ$, where $\frac{3}{5}<s\le1$.
Suppose also that every planar lattice projection of $Q$ has width
at least three and area at least $b$. Set
\[
 J=L(K),\qquad A=\Phi(\frac{9}{5})=\frac{81}{40},\qquad q=\frac{5s}{5s-3}.
\]
If $J\le q$, then $|K|\le\max\{\frac{729}{128},Aq\}$. If $J\ge q$,
then $|K|\le AJ$. Furthermore, for every $\gamma>0$ such that
$A\gamma^2\le b$ and $\gamma s>1$, we have
\begin{equation}\label{aux:eq-hollow-envelope}
 |K|\le H_\gamma(s):=
 \max\left\{\frac{729}{128},\frac{5As}{5s-3},
                         \frac{A(\gamma s)^3}{3(\gamma s-1)}\right\}.
\end{equation}
For fixed $\gamma$, $H_\gamma$ is convex on $\frac{3}{5}<s\le1$, $\gamma s>1$.
\end{lemma}

\begin{proof}
We project along a direction attaining $J$. For $0\le t\le J$, let
$E_t$ be the set of points over which the fiber of $K$ has length at
least $t$. The image $E_0$ contains a translate of a planar projection
of $sQ$, so $|E_0|\ge bs^2$ and $\lw(E_0)\ge w:=3s$. If $J>1$,
\cref{aux:chords} implies $I(E_1)\le I(K)=0$. It also implies that
$E_1$ contains a translate of $\left(1-\frac{1}{J}\right)E_0$.
Consequently, $\lw(E_1)\ge v:=w\left(1-\frac{1}{J}\right)$.
The value of $v$ increases with $J$ and equals $\frac{9}{5}$ at $J=q$.
Since $s\le1$, we have $q\ge\frac{5}{2}$.

First suppose $J\le q$. If $J\le\frac{9}{4}$, we use
\cref{aux:torus} with $T=\frac{9}{4}$ and $j=2$. This is possible
because $\mu(K)\ge1$ for a hollow body. Thus
$|K|\le\frac{1}{2}\left(\frac{9}{4}\right)^3=\frac{729}{128}$.
For $\frac{9}{4}\le J\le q$, we have
$1<\frac{5w}{9}\le v\le\frac{9}{5}$, where the strict inequality
uses $w=3s>\frac{9}{5}$. By \cref{aux:planar},
$|E_1|\le\Phi(v)$. From \eqref{aux:support-basic}, with $a=1$
and $n=2$, the upper bound for $|K|$ is the larger of
\[
  J\Phi(v)=\frac{w^2(J-1)^2}{2((w-1)J-w)},\qquad
  \frac{J^3\Phi(v)}{3(J-1)^2}
        =\frac{w^2J^2}{6((w-1)J-w)}.
\]
Both denominators are positive because $v>1$. Each expression is
convex in $J$. In fact, if $p$ and $h$ are affine functions with
$h>0$, the second derivative of $\frac{p^2}{h}$ is
$\frac{2(p'h-ph')^2}{h^3}\ge0$. Thus the larger of the two
expressions is also convex, and we need only evaluate it at the
endpoints. Their ratio, second to first, is $\frac{J^2}{3(J-1)^2}$.
At $J=\frac{9}{4}$ it is $\frac{27}{25}>1$, so the larger value is
\[
       \frac{27w^2}{8(5w-9)}=\frac{w}{3}Aq\le Aq .
\]
Here $q=\frac{5w}{5w-9}$ and $w\le3$. For $J\ge q\ge\frac{5}{2}$,
the ratio is at most $\frac{25}{27}<1$. In particular the larger value
at $J=q$ is $q\Phi(\frac{9}{5})=Aq$. This proves the assertion for
$J\le q$. If $J\ge q$, then $v\ge\frac{9}{5}$, and \cref{aux:planar}
implies $|E_1|\le A$. The same ratio comparison in
\eqref{aux:support-basic} now yields $|K|\le AJ$.

For \eqref{aux:eq-hollow-envelope} it remains to consider $J\ge q$.
We use \cref{aux:support} with $n=2$ and $a=p=1$. Since
$|E_1|\le A$ and $|E_0|\ge bs^2$, its supporting line has a
parameter $c$ satisfying
\[
 0<\gamma s-1\le\sqrt{\frac{bs^2}{A}}-1
 \le c\le\frac{1}{J-1}\le\frac{2}{3},\qquad
 |K|\le A\frac{(1+c)^3}{3c}.
\]
The volume estimate here is \eqref{eq:moment-cone}. The function
$\frac{(1+c)^3}{3c}=\frac{1}{3c}+1+c+\frac{c^2}{3}$ is convex for
$c>0$. At $c=\gamma s-1$ the bound is the last term in
\eqref{aux:eq-hollow-envelope}, and at $c=\frac{2}{3}$ it is
$\frac{75}{16}<\frac{729}{128}$. This proves the required maximum.
Finally, $\frac{5s}{5s-3}=1+\frac{3}{5s-3}$ and
$\frac{x^3}{x-1}=x^2+x+1+\frac{1}{x-1}$ are convex for
$s>\frac{3}{5}$ and $x>1$, respectively. Hence $H_\gamma$ is convex
on its stated domain.
\end{proof}

\subsection{Areas of planar fibers}

For a planar projection of a four-dimensional body, the square root
of the fiber area is concave. The following estimate for its second
moment will apply both to the original body and to a translate of its
dilate by $\sigma_0=\frac{99}{100}$. We use this fixed value of $\sigma_0$
throughout this section.

\begin{lemma}\label{aux:planar-moment}
Let $B$ be a planar convex body with
$\lw(B)\ge3\sigma_0$ and $|B|\ge\frac{7}{2}\sigma_0^2$. If
$f\colon B\to[0,\infty)$ is upper semicontinuous and concave, and
$f(z)\le1$ for each interior lattice point $z$ of $B$, then
\[
                     \int_B f^2\le\max\{\frac{25|B|}{9},\frac{58}{5}\}.
\]
\end{lemma}

\begin{proof}
Let $M=\max_B f$. If $M\le\frac{5}{3}$, then $\int_B f^2\le\frac{25|B|}{9}$.
Suppose that $M>\frac{5}{3}$. Let $E=\{y\in B:f(y)\ge1\}$ and $w=3\sigma_0$. For a point
$y_*\in B$ attaining $M$, concavity implies
\[
                 (1-\frac{1}{M})B+\frac{y_*}{M}\subset E.
\]
Indeed, the value of $f$ at every point of the displayed translate
is at least $(1-\frac{1}{M})\,0+\frac{1}{M}M=1$. For
$z\in\intr E$ with $z\ne y_*$, the segment from $y_*$ through $z$
extends beyond $z$ in $E$. The value of $f$ at $y_*$ exceeds one and
at the far endpoint it is at least one, so $f(z)>1$ by concavity.
This also holds at $y_*$ itself. Since $\intr E\subseteq\intr B$
and $f\le1$ at the interior lattice points of $B$, $E$ is hollow. The displayed inclusion also gives
$\lw(E)\ge w(1-\frac{1}{M})$.

By \cref{aux:support} at level one with $n=p=2$, there exists
$\max\{0,\sqrt{\frac{|B|}{|E|}}-1\}\le c\le\frac{1}{M-1}$ such that
\[
                    \int_B f^2\le |E|\int_0^M 2t(1+c(1-t))^2\,dt .
\]
The integral on the right is convex in $c$, as its integrand is a
nonnegative multiple of the square of an affine function of $c$.
At $c=0$ the integral is $M^2$. At $c=\frac{1}{M-1}$ it is
\[
 \frac{1}{(M-1)^2}\int_0^M 2t(M-t)^2\,dt
       =\frac{M^4}{6(M-1)^2}.
\]
Thus, on the full interval $0\le c\le\frac{1}{M-1}$,
\begin{equation}\label{eq:planar-moment-endpoints}
   \int_B f^2\le |E|\max\left\{M^2,\frac{M^4}{6(M-1)^2}\right\}.
\end{equation}

Suppose first that $\frac{5}{3}\le M\le\frac{7}{3}$. The lower
bound $w(1-\frac{1}{M})$ on $\lw(E)$ varies between
$\frac{2w}{5}=\frac{297}{250}>1$ and $\frac{4w}{7}=\frac{297}{175}<2$.
Hence $|E|\le\Phi(w(1-\frac{1}{M}))$ by \cref{aux:planar}. After
substitution in \eqref{eq:planar-moment-endpoints}, the two quantities
to be bounded are
\[
 \frac{w^2 M(M-1)^2}{2((w-1)M-w)}
 \qquad\text{and}\qquad
 \frac{w^2 M^3}{12((w-1)M-w)} .
\]
The pole $\frac{w}{w-1}$ lies strictly between $1$ and $\frac{5}{3}$.
When either numerator is divided by the affine denominator, the
quotient is a quadratic with positive leading coefficient, and the
remainder is the numerator evaluated at the pole, hence positive.
The reciprocal of the denominator is convex to the right of the pole.
Both quantities are therefore convex in $M$. The two endpoint
comparisons reduce to
\[
\begin{array}{c|cc}
 M&\Phi(w(1-\frac{1}{M}))&
       \max\{M^2,\frac{M^4}{6(M-1)^2}\}\\ \hline
 \frac{5}{3}&\frac{88209}{23500}<\frac{19}{5}&\frac{625}{216}\\[2pt]
 \frac{7}{3}&\frac{88209}{42700}<\frac{207}{100}&\frac{49}{9}
\end{array}
\]
The products in the two rows are less than $11$ and $\frac{23}{2}$,
respectively. Both are less than $\frac{58}{5}$.

If $M\ge\frac{7}{3}$, then $\lw(E)\ge\frac{4w}{7}$, so the same
planar bound yields $|E|<\frac{207}{100}$. Since
$(\frac{9}{7})^2\frac{207}{100}<\frac{7}{2}\sigma_0^2\le |B|$,
we have $\sqrt{\frac{|B|}{|E|}}-1>\frac{2}{7}$. The upper bound
$c\le\frac{1}{M-1}$ is at most $\frac{3}{4}$. The logarithmic
derivative of $\frac{(1+c)^4}{c^2}$ is
$\frac{2(c-1)}{c(1+c)}\le0$ on $(0,1]$. Thus \eqref{eq:moment-cone},
followed by the lower bound on $c$, implies
\[
 \int_B f^2\le |E|\frac{(1+c)^4}{6c^2}
 \le\frac{207}{100}\frac{(\frac{9}{7})^4}{6(\frac{2}{7})^2}
 =\frac{452709}{39200}<\frac{58}{5}. \qedhere
\]
\end{proof}

\begin{lemma}
Let $W$ be a four-dimensional convex body, and let $h$ be a positive
integer with $I(W)<h$. Set $x=\frac{L(W)}{h}$ and
\[
 G(x)=
 \begin{cases}
       \frac{9}{4},&0<x\le\frac{3}{2},\\
       \frac{2x^2}{2x-1},&x\ge\frac{3}{2} .
 \end{cases}
\]
Let $B$ be a planar lattice projection of $W$. For each interior
lattice point $z$ of $B$, the fiber $W_z$ has area at most $hG(x)$.
If also $\lw(B)\ge3\sigma_0$ and $|B|\ge\frac{7}{2}\sigma_0^2$, then
\begin{equation}\label{eq:four-fiber-area}
             \frac{|W|}{h}\le G(x)\max\{\frac{25|B|}{9},\frac{58}{5}\}.
\end{equation}
\end{lemma}

\begin{proof}
For an interior lattice point $z$ of $B$, \eqref{eq:relative-fiber}
identifies the relative interior lattice points of $W_z$ with the
interior lattice points of $W$ over $z$. Thus $I(W_z)\le h-1$.
The fiber carries an affine translate of the kernel lattice.
We choose lattice coordinates on it and a primitive height attaining
$u=\lw(W_z)$. The kernel of this height in the fiber lattice is
generated by a vector primitive also in the ambient lattice. Hence
a chord on a line of constant height has lattice length at most
$L(W)$. Since the primitive height ranges over an interval of length
$u$, Fubini's theorem yields $|W_z|\le uL(W)$.

For $u\le1$, this proves $\frac{|W_z|}{h}\le x$. For $u\ge2$,
\cref{aux:planar} bounds $|W_z|$ by $2$ if $I(W_z)=0$.
Otherwise $h\ge2$ and the same lemma bounds it by
$2(h-1)+\frac{5}{2}=2h+\frac{1}{2}\le\frac{9h}{4}$.
Thus $\frac{|W_z|}{h}\le\frac{9}{4}$ in either case. If $1<u\le2$,
then $I(W_z)+1\le h$, and \cref{aux:planar} implies
\[
                 \frac{|W_z|}{h}\le\min\{xu,\Phi(u)\}.
\]
For $x\le1$ the first term is at most two. For $x\ge1$, it increases
with $u$, whereas $\Phi(u)$ decreases on $(1,2]$. The two are equal
at $u=\frac{2x}{2x-1}\in(1,2]$. Below this value the first term is at
most its value at the crossing, and above it the second is at most its
value there. The resulting bound is $\frac{2x^2}{2x-1}$.
It exceeds $x$, has derivative $\frac{4x(x-1)}{(2x-1)^2}\ge0$ for
$x\ge1$, and equals $\frac{9}{4}$ at $x=\frac{3}{2}$. This proves
$|W_z|\le hG(x)$ in all three width ranges.

For the volume inequality, the function
$f(y)=\sqrt{\frac{|W_y|}{hG(x)}}$ is upper semicontinuous and concave
by the fiber properties in \cref{sec:reduction}. We have just proved
that $f(z)\le1$ for every interior lattice point of $B$. From
\cref{aux:planar-moment} and $\int_B |W_y|\,dy=|W|$ we obtain
\eqref{eq:four-fiber-area}.
\end{proof}

\subsection{Hollow sets of long chords}

The inequalities $71<A_4(1)=\frac{147}{2}$ and $\frac{71}{\sigma_0^4}<74<V_4(1)$
explain the common constant in the next lemma. We apply it to a
hollow translate of a dilated four-dimensional image in
\cref{aux:four-profile} and to the lattice four-polytope at
the end of the section.

\begin{lemma}\label{lem:hollow-chord-four}
Let $W$ be a four-dimensional convex body. For a fixed primitive
lattice direction, let $L$ be the maximum chord length in that
direction, and let $D_t$ be the set of projected points whose
chord has length at least $t$. Suppose that $D_a$ is hollow for some
$0<a<L$ and that $D_0$ contains a translate of $\sigma Q$, where
$\sigma_0\le\sigma\le1$. Assume that $Q$ is a three-dimensional convex
body of volume at least $\frac{43}{6}$ and that each of its planar lattice
projections has width at least three and area at least $b\ge\frac{7}{2}$.
Then $|W|<71a$ whenever one of the following conditions holds.
\begin{enumerate}
\item $\sigma(1-\frac{a}{L})\ge\frac{81}{100}$.
\item $\sigma=1$, $L\ge3a$, and $b\ge5$.
\item $\sigma=\sigma_0$, $L\ge\frac{41a}{10}$, and $b\ge4$.
\end{enumerate}
\end{lemma}

\begin{proof}
In case (i), $\sigma\le1$ implies
$\frac{a}{L}\le1-\frac{81}{100}=\frac{19}{100}$. In (ii) and (iii)
the stated bounds on $L$ also imply $L\ge3a$. Thus
$4(1-\frac{a}{L})^3\ge4(\frac{2}{3})^3>1$. By \eqref{aux:support-basic}
with $n=3$, the larger of its two terms is $L$, so $|W|\le L|D_a|$.
The inclusion of \cref{aux:chords} implies that $D_a$ contains a
translate of $\alpha Q$, where $\alpha=\sigma(1-\frac{a}{L})$.
In all three cases $\alpha>\frac{3}{5}$. Since
$\frac{a}{L}=1-\frac{\alpha}{\sigma}$, \cref{aux:hollow-envelope}
therefore implies
\[
 \frac{|W|}{a}\le\frac{H_\gamma(\alpha)}{1-\frac{\alpha}{\sigma}}
 \quad\text{if }\frac{81}{40}\gamma^2\le b\text{ and }\gamma\alpha>1.
\]

Suppose first that $\alpha\le\frac{9}{10}$. For
$\alpha\ge\frac{81}{100}$ we use $\gamma=\frac{21}{16}$ in all
three cases. Since $\sigma\ge\sigma_0$, it suffices on this interval
to prove $H_\gamma(\alpha)<71(1-\frac{\alpha}{\sigma_0})$.
For $\alpha<\frac{81}{100}$, only cases (ii) and (iii) remain.
Their lower endpoints for $\alpha$ are, respectively,
$\frac{2}{3}$ and $\sigma_0(1-\frac{10}{41})=\frac{3069}{4100}$.
We use $\gamma=\frac{31}{20}$ in (ii) and $\gamma=\frac{7}{5}$ in
(iii). The three choices satisfy the area condition since
\[
 \frac{81}{40}\left(\frac{21}{16}\right)^2<\frac{7}{2},\qquad
 \frac{81}{40}\left(\frac{31}{20}\right)^2<5,\qquad
 \frac{81}{40}\left(\frac{7}{5}\right)^2<4 .
\]
On each interval $H_\gamma$ is convex, while
$71(1-\frac{\alpha}{\sigma})$ is affine for fixed $\sigma$. It is
therefore enough to establish the strict inequality at both endpoints
of each interval.

We now compare the three terms at the endpoints in
\eqref{aux:eq-hollow-envelope}. In the table the first pair has
$\sigma=1$, and the other two pairs use $\sigma=\sigma_0$. For the
third term write $v=\gamma\alpha$. The last column is the upper bound
we seek at the given endpoint.
\[
\renewcommand{\arraystretch}{1.35}
\begin{array}{cc|ccc|c}
 \gamma&\alpha& v&
 \frac{405\alpha}{40(5\alpha-3)}\text{ at most}&
 \frac{1}{v-1}\text{ at most}& 71(1-\frac{\alpha}{\sigma})\\ \hline
 \frac{31}{20}&\frac{2}{3}&\frac{31}{30}&\frac{81}{4}&30&\frac{71}{3}\\
 \frac{31}{20}&\frac{81}{100}&\frac{2511}{2000}&8&4&\frac{1349}{100}\\
 \frac{7}{5}&\frac{3069}{4100}&\frac{21483}{20500}&11&21&\frac{710}{41}\\
 \frac{7}{5}&\frac{81}{100}&\frac{567}{500}&8&8&\frac{142}{11}\\
 \frac{21}{16}&\frac{81}{100}&\frac{1701}{1600}&8&\frac{317}{20}&\frac{142}{11}\\
 \frac{21}{16}&\frac{9}{10}&\frac{189}{160}&\frac{243}{40}&\frac{50}{9}&\frac{71}{11}
\end{array}
\]
In particular $v>1$ at each left endpoint, so $\gamma\alpha>1$ on
all three intervals. The constant term $\frac{729}{128}$ is less
than the least entry $\frac{71}{11}$ in the last column. The fourth
column bounds the second term strictly below the last column in each
row. For the third term we use the identity
\[
       \frac{81}{40}\frac{v^3}{3(v-1)}
       =\frac{27}{40}\left(v^2+v+1+\frac{1}{v-1}\right).
\]
All six values of $v$ are less than $\frac{13}{10}$, so
$v^2+v+1<4$. In the fifth row $v<\frac{27}{25}$, which implies the
sharper bound $v^2+v+1<\frac{13}{4}$. Thus the third terms, in row
order, are bounded above by $\frac{27}{40}$ times
\[
 34,\quad 8,\quad 25,\quad 12,\quad
 \frac{13}{4}+\frac{317}{20},\quad 4+\frac{50}{9}.
\]
These products are respectively less than the entries in the last
column. In particular the last two are $\frac{5157}{400}<\frac{142}{11}$
and $\frac{129}{20}<\frac{71}{11}$. The endpoint comparisons imply
$\frac{|W|}{a}<71$ on each interval.

Suppose next that $\alpha\ge\frac{9}{10}$. A translate of
$\frac{9}{10}Q$ then lies in $D_a$. In the last row,
$v=\frac{189}{160}<\frac{19}{16}$ implies $v^2+v+1<\frac{18}{5}$.
Thus the third term of $H_\gamma(\frac{9}{10})$ is at most
$\frac{27}{40}(\frac{18}{5}+\frac{50}{9})=\frac{309}{50}$.
From \eqref{aux:eq-hollow-envelope} with $\gamma=\frac{21}{16}$
we obtain
\[
 |D_a|\le H_\gamma\left(\frac{9}{10}\right)
 \le\max\left\{\frac{729}{128},\frac{243}{40},
              \frac{309}{50}\right\}<\frac{25}{4}.
\]
Also $|D_0|\ge\sigma_0^3\frac{43}{6}>\frac{25}{4}(\frac{31}{30})^3$
and $\frac{a}{L}=1-\frac{\alpha}{\sigma}\le\frac{1}{10}$. Thus
$\sqrt[3]{\frac{4|D_0|}{25}}-1>\frac{1}{30}$ and
$\frac{L}{a}\ge10>4$. These are the hypotheses of \eqref{aux:cone}
with $n=3$, $A=\frac{25}{4}$ and $\eta=\frac{1}{30}$. It follows that
\[
 \frac{|W|}{a}\le\frac{25}{4}\left(\frac{31}{30}\right)^4
                          \frac{30}{4}
              =\frac{923521}{17280}<54<71. \qedhere
\]
\end{proof}

\begin{proposition}\label{aux:four-profile}
Let $V$ be a four-dimensional convex body with $\lw(V)\ge3$.
Suppose that every planar lattice projection of $V$ has area at least
$\frac{7}{2}$ and every three-dimensional lattice projection has volume at least
$\frac{43}{6}$. If $\mu(V)>\frac{99}{100}$, then $|V|<74$.
\end{proposition}

\begin{proof}
Since $\sigma_0<\mu(V)$, there is a hollow translate $W$ of
$\sigma_0V$. Its volume is $\sigma_0^4|V|$. Suppose for a
contradiction that $|W|\ge74\sigma_0^4$, and set $L=L(W)$.
If $L\le\frac{41}{10}$, then $4<\frac{41}{10}\mu(W)$ since
$\mu(W)\ge1$. By \cref{aux:torus} with $T=\frac{41}{10}$ and $j=4$,
\[
 |W|\le\frac{41^4}{4\cdot10^4}<71
       <74\left(\frac{99}{100}\right)^4 .
\]
Hence $L>\frac{41}{10}$.

Suppose $L\le\frac{11}{2}$. If some planar image of $V$ had area
at most four, the corresponding image of $W$ would have width at least
$3\sigma_0$ and area between $\frac{7}{2}\sigma_0^2$ and
$4\sigma_0^2$. We can use \eqref{eq:four-fiber-area} with $h=1$,
since $W$ is hollow. The monotonicity of $G$ and the inequality
$\frac{25}{9}\,4\sigma_0^2<\frac{58}{5}$ would imply
\[
 |W|\le G\left(\frac{11}{2}\right)\frac{58}{5}
       =\frac{121}{20}\frac{58}{5}
       =\frac{3509}{50}<71<74\sigma_0^4 .
\]
Thus every planar image of $V$ has area greater than four in this
range of $L$.

We now project $V$ and $W$ along a direction attaining $L(W)$.
If $Q$ is the image of $V$, then the image of $W$ is a translate of
$\sigma_0Q$. Every planar lattice projection of $Q$ is a planar
lattice projection of $V$, and $|Q|\ge\frac{43}{6}$ by hypothesis.
Since $I(W)=0$ and $1<L$, the chord superlevel set $D_1$ of $W$ is
hollow by \cref{aux:chords}. For $\frac{41}{10}<L\le\frac{11}{2}$,
we can take $b=4$ in \cref{lem:hollow-chord-four}(iii).
For $L\ge\frac{11}{2}$ we have
$\sigma_0(1-\frac{1}{L})\ge\frac{99}{100}\frac{9}{11}
=\frac{81}{100}$, so (i) of the same lemma applies. In either case,
with $a=1$ we obtain $|W|<71<74\sigma_0^4$, the desired contradiction.
\end{proof}

\subsection{Lattice-point counts for shorter chords}

\begin{lemma}\label{lem:short-layer}
Let $q\ge2$ be a power of two, and let $C$ be a three-dimensional
convex body with $I(C)<q$. Suppose that $C$ contains a translate
of $sQ$, where $Q$ is a three-dimensional convex body whose planar
lattice projections have width at least three and area at least $b$.
Then
\[
 \begin{aligned}
 |C|&\le \frac{45q}{4} &&\text{if }s=\frac{2}{3},\ b\ge\frac{15}{2},\\
 |C|&\le 14q   &&\text{if }s=\frac{7}{12},\ b\ge10 .
 \end{aligned}
\]
\end{lemma}

\begin{proof}
Let $J=L(C)$, and let $E_m$ be the set of points over which the
fiber has length at least $m$ in a direction attaining $J$. The image
$E_0$ contains a translate of $sB$, where $B$ is the projection of $Q$
along the same direction. By \eqref{aux:planar-radius} with width
lower bound three,
\[
 \min\{\mu(B),\frac{2}{3}\}\le\frac{1}{3}+\frac{3}{2|B|}
                         \le\frac{1}{3}+\frac{3}{2b}.
\]
The right side is less than $\frac{2}{3}$ in both cases, so it
bounds $\mu(B)$. If $J>q$, then $qI(E_q)\le I(C)<q$
by \cref{aux:chords}. Thus $E_q$ is hollow and contains a translate
of $s(1-\frac{q}{J})B$, whence $s(1-\frac{q}{J})\le\mu(B)$.
This inequality also holds for $J\le q$, when its left side is nonpositive.

If $s=\frac{2}{3}$ and $b\ge\frac{15}{2}$, then
$\mu(B)\le\frac{1}{3}+\frac{1}{5}=\frac{8}{15}$. Hence
$\frac{q}{J}\ge1-\frac{8}{15}\frac{3}{2}=\frac{1}{5}$ and $J\le5q$.
If $s=\frac{7}{12}$ and $b\ge10$, we obtain
$\mu(B)\le\frac{1}{3}+\frac{3}{20}=\frac{29}{60}$ and
$\frac{q}{J}\ge1-\frac{29}{60}\frac{12}{7}=\frac{6}{35}$.
We set $a=\frac{5}{2}$ in the first case and $a=3$ in the second.
Thus $J\le2aq$ in both. If $J<a$, then $|C|<a^3$ by \cref{aux:torus}, and the result follows from
$(\frac{5}{2})^3<\frac{45q}{4}$ and $3^3<14q$ for $q\ge2$.

For $a\le J\le2aq$, the intervals $[am,2am]$ with $m=1,2,4,\ldots,q$
cover $[a,2aq]$, since successive intervals share an endpoint.
Thus some divisor $m$ of $q$ satisfies $a\le t:=\frac{J}{m}\le2a$.
Set $w=3s$. The planar image of $C$ has width at least $w$, and
the lower endpoint of $w(1-\frac{1}{t})$ is $\frac{6}{5}$ in the
first case and $\frac{7}{6}$ in the second. In each case
$1<w(1-\frac{1}{t})\le w\le2$ and $t\ge\frac{5}{2}$, so the
divisibility assertion of \cref{lem:three-count} yields
\[
 \frac{|C|}{q}\le t\Phi_+\left(w\left(1-\frac{1}{t}\right)\right)
   =\max\left\{\frac{9t}{4},\frac{w^2(t-1)^2}{2((w-1)t-w)}\right\}.
\]
The denominator is positive because $w(1-\frac{1}{t})>1$. The
second term is convex as the square of an affine function divided by
a positive affine function. The first term is affine, so the maximum
is convex and its largest value occurs at an endpoint. The four
substitutions are
\[
\renewcommand{\arraystretch}{1.3}
\begin{array}{cc|cc}
 w&t&\frac{9t}{4}&\frac{w^2(t-1)^2}{2((w-1)t-w)}\\ \hline
 2&\frac{5}{2}&\frac{45}{8}&9\\
 2&5&\frac{45}{4}&\frac{32}{3}\\
 \frac{7}{4}&3&\frac{27}{4}&\frac{49}{4}\\
 \frac{7}{4}&6&\frac{27}{2}&\frac{1225}{88}
\end{array}
\]
The first two rows are bounded by $\frac{45}{4}$, and the last two
by $14$. This proves both estimates.
\end{proof}

\begin{proof}[Proof of \cref{prop:projections} for $(d,u)=(4,1)$]
We argue by contradiction. Suppose that all lattice projections of the
given four-polytope $P$ of ranks two and three exceed $V_2(1)=3$ and
$V_3(1)=7$, respectively. Their volumes are then at least
$\frac{7}{2}$ and $\frac{43}{6}$. Let $K=k+1$, $L=L(P)$, and
$X=\frac{L}{K}$. By hypothesis $\frac{|P|}{K}\ge A_4(1)=\frac{147}{2}$.
Every planar lattice projection $B$ has width at least three, and
its area $b=|B|$ satisfies $2b\in\Z$. Thus we may apply
\eqref{eq:four-fiber-area} with $h=K$ to obtain
\[
       \frac{147}{2}\le G(X)\max\left\{\frac{25b}{9},\frac{58}{5}\right\}.
\]
The right side is nondecreasing in both $X$ and $b$. At the following
three pairs, direct substitution gives
\[
\renewcommand{\arraystretch}{1.3}
\begin{array}{cc|ccc}
 X&b&G(X)&\max\{\frac{25b}{9},\frac{58}{5}\}&
                G(X)\max\{\frac{25b}{9},\frac{58}{5}\}\\ \hline
 \frac{100}{19}&\frac{9}{2}&\frac{20000}{3439}&
              \frac{25}{2}&\frac{250000}{3439}\\
 3&7&\frac{18}{5}&\frac{175}{9}&70\\
 2&\frac{19}{2}&\frac{8}{3}&\frac{475}{18}&\frac{1900}{27}
\end{array}
\]
Each entry in the last column is less than $\frac{147}{2}$. For example,
if $X\le3$, the second row rules out $b\le7$, and $2b\in\Z$
then implies $b\ge\frac{15}{2}$. The same argument in the other two
rows establishes
\begin{equation}\label{eq:four-area}
 b\ge5\quad (X\le\frac{100}{19}),\qquad
 b\ge\frac{15}{2}\quad (X\le3),\qquad b\ge10\quad (X\le2).
\end{equation}

We project along a direction attaining $L$. Write $Q$ for the image
and $D_t$ for the set of projected points whose fibers have length
at least $t$. Suppose first that $X\ge3$. Then $K<L$ and
$K I(D_K)\le I(P)=K-1$, so $D_K$ is hollow. Every planar lattice
projection of $Q$ is also a planar lattice projection of $P$.
For $3\le X\le\frac{100}{19}$, \eqref{eq:four-area} allows
$b=5$ in \cref{lem:hollow-chord-four}(ii) with $a=K$ and $\sigma=1$.
For $X\ge\frac{100}{19}$ we have $1-\frac{1}{X}\ge\frac{81}{100}$,
so (i) of the same lemma applies. In either case
$\frac{|P|}{K}<71<\frac{147}{2}$, a contradiction.

Now suppose $X<3$. Let $q\ge2$ be the smallest power of two such that
$qX\ge\frac{16}{5}$, and set $n=\lceil\frac{K}{q}\rceil$. If $q>2$,
then $\frac{q}{2}X<\frac{16}{5}$ by minimality. If $q=2$, then
$qX<6<\frac{32}{5}$. Thus
$\frac{16}{5}\le qX\le\frac{32}{5}$ in both cases. Whenever $n<L$,
\cref{aux:chords} shows that $nI(D_n)\le K-1<nq$, since $K\le nq$.
In particular $I(D_n)<q$, and $D_n$ contains a translate of
$(1-\frac{n}{L})Q$. By \eqref{aux:support-basic} with $a=n$, the
volume of $P$ is at most the larger of $L|D_n|$ and
$\frac{L^4|D_n|}{4(L-n)^3}$. Equivalently,
\begin{equation}\label{eq:four-chord-volume}
 \frac{|P|}{K}\le X|D_n|
   \max\left\{1,\frac{1}{4(1-\frac{n}{L})^3}\right\}.
\end{equation}

If $n\le\frac{L}{3}$, we have $1-\frac{n}{L}\ge\frac{2}{3}$.
In particular $n<L$, and $4(\frac{2}{3})^3>1$ makes the maximum
in \eqref{eq:four-chord-volume} equal to one. By convexity, a
translate of $\frac{2}{3}Q$ lies in $D_n$. Since $X<3$,
\eqref{eq:four-area} bounds every planar image area below by
$\frac{15}{2}$. From \cref{lem:short-layer},
\[
         \frac{|P|}{K}\le X|D_n|\le \frac{45}{4}qX
               \le \frac{45}{4}\frac{32}{5}=72<\frac{147}{2}.
\]

Suppose instead that $n>\frac{L}{3}$ and $qn\le\frac{4K}{3}$.
Then $qX=\frac{qL}{K}<\frac{3qn}{K}\le4$, so $X<2$ since $q\ge2$.
Moreover
\[
 \frac{n}{L}=\frac{qn}{K}\frac{1}{qX}
       \le\frac{4}{3}\frac{5}{16}=\frac{5}{12}.
\]
Thus $n<L$ and $D_n$ contains a translate of $\frac{7}{12}Q$.
Now \eqref{eq:four-area} allows $b=10$ in \cref{lem:short-layer},
whence $|D_n|\le14q$. The maximum in \eqref{eq:four-chord-volume}
is at most $\frac{12^3}{4\cdot7^3}=\frac{432}{343}>1$. Therefore
\[
        \frac{|P|}{K}\le14qX\frac{432}{343}
          <56\frac{432}{343}=\frac{3456}{49}<\frac{147}{2}.
\]

The remaining case is $n>\frac{L}{3}$ and $qn>\frac{4K}{3}$.
Since $n-1<\frac{K}{q}<\frac{3n}{4}$, we have $n<4$. On the
other hand, $K\ge2$ and \cref{aux:torus} imply
$L^4\ge|P|\ge147>3^4$. Thus $n>\frac{L}{3}>1$, so $n=2$ or $n=3$.
From \eqref{eq:gzpw-density} and $k=K-1$ we obtain
\[
 1-\mu(P)\le\left(\frac{K-1}{|P|}\right)^{\frac{1}{4}}
       <\left(\frac{2}{147}\right)^{\frac{1}{4}}<\frac{7}{20}.
\]
The last comparison follows from $2\cdot20^4=320000
<352947=147\cdot7^4$. We can therefore take $T=\frac{40n}{13}$ and
$j=2n$ in \cref{aux:torus}. Indeed, $T>3n>L$ and
$\mu(P)T>\frac{13}{20}T=2n$. Consequently
\[
 |P|\le\frac{1}{2n}\left(\frac{40n}{13}\right)^4
     =\frac{1280000}{28561}n^3<45n^3 .
\]
For $n=2$ we have $q<K<\frac{3q}{2}$. If $q=2$ this would require
$2<K<3$, so $q\ge4$ and $K\ge5$. For $n=3$ we have
$2q<K<\frac{9q}{4}$. At $q=2$ and $q=4$ these open intervals are
$(4,\frac{9}{2})$ and $(8,9)$, neither of which contains an integer.
Hence $q\ge8$ and $K\ge17$. The two lower bounds on $K$ now yield
\[
 \frac{|P|}{K}<
 \max\left\{\frac{45\cdot2^3}{5},\frac{45\cdot3^3}{17}\right\}
 =72<\frac{147}{2}.
\]
This last contradiction completes the proof.
\end{proof}

\section{Reduction to a four-dimensional lattice image}\label{sec:gzpw-wide}

We assume the hypotheses of \cref{prop:projections} with $u=1$ and
$d\ge5$. Write $V_i=V_i(1)$ and $A_i=A_i(1)$. Suppose that every
proper lattice projection of rank $i\ge2$ has volume greater than
$V_i$. We use the successive projections of \cref{aux:flag}, with
$Q_d=C$, $L_i=L(Q_i)$ and $\mu_i=\mu(Q_i)$. We have $|Q_i|>V_i$
for $2\le i\le d$, since $|Q_d|\ge2A_d>V_d$. In particular,
$|Q_4|>\frac{1806}{24}>75$. All lattice projections have width at
least three. By our assumption on all proper lattice projections, each planar image has
area greater than $V_2=3$ and each three-dimensional image has volume
greater than $V_3=7$. Their normalized volumes are integers. Hence the
respective lower bounds are $\frac{7}{2}$ and $\frac{43}{6}$. The same argument applies to every planar and
three-dimensional projection of $Q_4$, by composition with the
projection onto $Q_4$. Thus \cref{aux:four-profile} and $|Q_4|>75$
imply $\mu_4\le0.99$. Our aim is to deduce the opposite strict
inequality from the covering radii of the successive images.
All terminating decimals denote exact rational numbers.

\subsection{Covering radii in dimensions four and five}

\begin{lemma}\label{lem:wide-radius}
Under the assumptions above,
\begin{equation}\label{eq:wide-loss-direct}
 \mu_5>0.87,\qquad \mu_4>0.81,\qquad
                 1-\mu_4<0.002+\frac{|Q_4|}{25000}.
\end{equation}
If $d\ge6$, then $\mu_5>0.95$.
\end{lemma}

\begin{proof}
For $u=1$ we have $V_i>A_i$, since
$f_i(1)=f_{i-1}(1)(f_{i-1}(1)+1)>f_{i-1}(1)^2$.
The recurrence gives $f_4(1)=1806$ and $f_5(1)=3263442$. Hence
\[
 \begin{aligned}
 A_5&=\frac{1806^2}{120}>27000,\\
 A_6&=\frac{3263442^2}{720}>14\cdot10^9>48^6,\\
 A_7&=\frac{(3263442\cdot3263443)^2}{5040}>1500^7 .
 \end{aligned}
\]
By \eqref{eq:volume-scale-growth}, $\sqrt[i]{A_i}$ and $\sqrt[i]{V_i}$
at least double at each step from $i=7$ onward.

Suppose first that $d=5$. Since $|Q_5|\ge(k+1)A_5$, we obtain from
\eqref{eq:gzpw-density}
\[
 1-\mu_5\le\sqrt[5]{\frac{k}{|Q_5|}}
 <\sqrt[5]{\frac{k}{27000(k+1)}}<0.13.
\]
Here $27000(0.13)^5>1$. Also $L_5^5\ge|Q_5|>54000>(8.4)^5$, so
$0.12L_5>1$. In particular,
$\lfloor0.87L_5\rfloor\ge0.87L_5-1>0.75L_5$.
This positive integer is at most $\mu_5 L_5$. By \cref{aux:torus}
with $T=L_5$ and $j=\lfloor0.87L_5\rfloor$, we have
$L_5^5\ge j|Q_5|>0.75L_5\cdot27000(k+1)$. Thus
$L_5^4>20250(k+1)$, and $\mu_4\ge\mu_5-\frac{1}{L_5}$ implies
\[
 1-\mu_4<
 \sqrt[5]{\frac{k}{27000(k+1)}}+
 \sqrt[4]{\frac{1}{20250(k+1)}} .
\]
For every positive integer $k$, the sum of these roots is less than $0.19$. The entries are upper bounds for the two
radicands, followed by the powers to which we compare them.
\[
\begin{array}{c|c|c}
 k & \text{first radicand} & \text{second radicand}\\ \hline
 1 & \frac{1}{54000}<(0.115)^5 & \frac{1}{40500}<(0.075)^4\\[2pt]
 2 & \frac{2}{81000}<(0.12)^5  & \frac{1}{60750}<(0.07)^4\\[2pt]
 k\ge3 & \frac{1}{27000}<(0.13)^5 & \frac{1}{81000}<(0.06)^4
\end{array}
\]
The sum of the two decimals in each row is $0.19$, so $\mu_4>0.81$.
For the last inequality in \eqref{eq:wide-loss-direct},
\eqref{eq:gzpw-one-step} with $\tau=A_5$ implies
$1-\mu_4\le\frac{|Q_4|}{A_5}<\frac{|Q_4|}{27000}$.

Now suppose $d\ge6$. By \eqref{eq:gzpw-density} and $|C|\ge(k+1)A_d$,
we have $1-\mu_d<\frac{1}{\sqrt[d]{A_d}}$. Moreover,
$\frac{1}{L_i}<\frac{1}{\sqrt[i]{V_i}}$, and
$\mu_i-\mu_{i-1}\le\frac{1}{L_i}$ by \cref{aux:torus}. For $d=6$
these bounds imply $1-\mu_5<\frac{2}{48}$. For $d\ge7$, the doubling
of the roots gives
\[
 \begin{aligned}
 1-\mu_6
 &\le1-\mu_d+\sum_{i=7}^{d}\frac{1}{L_i}\\
 &<\frac{1}{1500}
 \left(\frac{1}{2^{d-7}}+\sum_{i=7}^{d}\frac{1}{2^{i-7}}\right)
 =\frac{2}{1500}.
 \end{aligned}
\]
Hence $1-\mu_5<\frac{2}{1500}+\frac{1}{48}$ in this case.
In both cases $\mu_5>0.95$. If $L_5\le8$, then $8\mu_5>7$
and \cref{aux:torus} with $T=8$, $j=7$ would yield
$|Q_5|\le\frac{8^5}{7}<27000$, a contradiction. Thus $L_5>8$, and
$\mu_4\ge\mu_5-\frac{1}{L_5}>0.95-\frac{1}{8}>0.81$.

It remains to bound $1-\mu_4$ in terms of $|Q_4|$. First,
$|Q_5|\le L_5|Q_4|$ implies
$\frac{1}{L_5}\le\frac{|Q_4|}{|Q_5|}<\frac{|Q_4|}{27000}$.
Also $L_5\le2L_6$ by \eqref{aux:nested} at level zero, so
\[
 |Q_6|\le L_6|Q_5|\le L_6 L_5|Q_4|\le2L_6^2|Q_4|.
\]
When $d=6$, we have $|Q_6|\ge2A_6$. Consequently,
$\frac{1}{L_6}\le\sqrt{\frac{|Q_4|}{A_6}}$, and
\eqref{eq:gzpw-continuous-count} with $m=4$ and $\tau=A_6$ yields
\[
 1-\mu_4\le2\sqrt{\frac{|Q_4|}{A_6}}+\frac{1}{L_5}+\frac{1}{L_6}
             \le\frac{1}{L_5}+3\sqrt{\frac{|Q_4|}{A_6}} .
\]
When $d\ge7$, the preceding bound on $1-\mu_6$ yields
\[
 1-\mu_4\le1-\mu_6+\frac{1}{L_6}+\frac{1}{L_5}
 <0.002+\sqrt{\frac{2|Q_4|}{|Q_6|}}+\frac{1}{L_5}.
\]
In either estimate the square-root term is at most
$3\sqrt{\frac{|Q_4|}{14\cdot10^9}}$, since $A_6,|Q_6|>14\cdot10^9$.
After division by $|Q_4|>75$, this term is less than
\[
 \frac{3}{\sqrt{75\cdot14\cdot10^9}}<\frac{1}{337500}
                 =\frac{1}{25000}-\frac{1}{27000},
\]
where the strict inequality follows by squaring.
Together with $\frac{1}{L_5}<\frac{|Q_4|}{27000}$, this proves
\eqref{eq:wide-loss-direct}.
\end{proof}

\subsection{The diameter and volume of the four-dimensional image}

\begin{proof}[Proof of \cref{prop:projections} for $u=1$ and $d\ge5$]
We shall prove first that $L_4<17$ and then that $|Q_4|<200$. The latter
bound implies $1-\mu_4<0.002+\frac{200}{25000}=0.01$ by
\cref{lem:wide-radius}, contrary to $\mu_4\le0.99$.

We begin with two bounds for $|Q_3|$. In \cref{aux:hollow-envelope},
take $\gamma=\frac{21}{16}$, where $A=\Phi(\frac{9}{5})=\frac{81}{40}$.
For $\frac{16}{21}<x\le1$, set
\[
 \Psi(x)=\frac{H_\gamma(x)}{x^3}
  =\max\left\{\frac{729}{128x^3},\frac{A}{x^2(x-0.6)},
                         \frac{A\gamma^3}{3(\gamma x-1)}\right\}.
\]
For $x\le\mu_3$ a translate of $xQ_3$ is hollow. Every planar image of
$Q_3$ has width at least three and area at least $\frac{7}{2}>A\gamma^2$.
Thus the lemma applies, and division by $x^3$ yields $|Q_3|\le\Psi(x)$.
We shall use the following strict upper bounds for the individual terms.
They follow by substitution, and each comparison has positive denominators.
\[
\begin{array}{c|ccc}
 x & \frac{729}{128x^3} & \frac{A}{x^2(x-0.6)}
   & \frac{A\gamma^3}{3(\gamma x-1)}\\ \hline
 0.80 & 12 & 16 & 31\\
 0.85 & 10 & 12 & 14\\
 0.93 & 7.1 & 7.1 & 7
\end{array}
\]
For example, the third term simplifies to $\frac{250047}{10240(21x-16)}$.
At $x=0.93$ the second term is
$\frac{81\cdot10^6}{40\cdot93^2\cdot33}<7.1$.
If $\mu_3\ge0.93$, the last row would imply $|Q_3|<7.1<\frac{43}{6}$.
We conclude that $\mu_3<0.93$.

Our second bound applies also for $x$ near $0.72$. The planar flatness
theorem of Hurkens~\cite{Hurkens} states that a hollow planar convex
body has lattice width at most $1+\frac{2}{\sqrt{3}}$ in the normalization
of \cite[Theorem~2.1]{AW}. A translate of $\mu_2Q_2$ is hollow, and
$\lw(Q_2)\ge3$, so
$\mu_2\le\frac{1}{3}(1+\frac{2}{\sqrt{3}})<0.72$.
By the projection inequality of \cref{aux:torus},
$\mu_3\le\mu_2+\frac{1}{L_3}$. For $0.72<x\le\mu_3$ we therefore have
$L_3<\frac{1}{x-0.72}$. A hollow translate of $xQ_3$ has rational
diameter $xL_3$. The first two bounds of \cref{aux:hollow-envelope}
apply on the two respective ranges of this diameter, so in all cases
\[
 x^3|Q_3|\le\max\left\{\frac{729}{128},\frac{Ax}{x-0.6},AxL_3\right\}
                   \le\frac{Ax}{x-0.72}.
\]
For the middle term this is immediate from $x-0.6>x-0.72$, and for
the last term it follows from the diameter bound. For the first term,
$x\mapsto\frac{x}{x-0.72}$ decreases and
$\frac{A(0.93)}{0.93-0.72}>\frac{729}{128}$. We have proved
\begin{equation}\label{eq:wide-flat-three}
 |Q_3|\le\frac{A}{x^2(x-0.72)}\qquad(0.72<x\le\mu_3).
\end{equation}

To bound $L_4$, we first exclude $L_4\ge150$. Since $L_4\le2L_5$ by
\eqref{aux:nested}, two applications of \cref{aux:torus} yield
\[
 \mu_3\ge\mu_5-\frac{1}{L_5}-\frac{1}{L_4}\ge\mu_5-\frac{3}{L_4}.
\]
For $d\ge6$, \cref{lem:wide-radius} implies $\mu_5>0.95$, whence
$L_4\ge150$ would imply $\mu_3>0.95-\frac{3}{150}=0.93$.
For $d=5$, the weaker bound $\mu_5>0.87$ first implies $\mu_3>0.85$.
The row $x=0.85$ in the table then implies $|Q_3|<14$. Since $A_5>27000$,
\eqref{eq:gzpw-continuous-count} with $m=3$, $\tau=A_5$ gives
\[
 1-\mu_3\le2\sqrt{\frac{|Q_3|}{A_5}}+\frac{1}{L_4}+\frac{1}{L_5}
           <2\sqrt{\frac{14}{27000}}+\frac{3}{150}<0.07 .
\]
The last comparison follows from $14<27000(0.025)^2$. Both cases
contradict $\mu_3<0.93$, so $L_4<150$.

By \eqref{eq:wide-loss-direct}, $0.998<\mu_4+\frac{|Q_4|}{25000}$.
Now $\mu_4\le\mu_3+\frac{1}{L_4}$ and $|Q_4|\le L_4|Q_3|$, hence
\begin{equation}\label{eq:wide-rectangle}
             0.998<\mu_3+\frac{1}{L_4}+\frac{L_4|Q_3|}{25000}.
\end{equation}
For $17\le L_4\le45$ we have
$\mu_3\ge\mu_4-\frac{1}{L_4}>0.81-\frac{1}{17}>0.75$.
For $45\le L_4\le150$, we have
$\mu_3\ge\mu_5-\frac{3}{L_4}>0.87-\frac{3}{45}>0.80$.
We substitute the bound \eqref{eq:wide-flat-three} at $x=\mu_3$ on
the first range and $\Psi(\mu_3)$ on the second range. Both bounds
are convex functions of $x$ on the intervals under consideration.
Indeed, each term has the form of a positive constant times a product
of negative powers of positive affine functions. Its logarithm is convex,
since $-\log$ is convex, and the exponential preserves convexity.
The maximum of the three terms defining $\Psi$ is therefore also convex.
For fixed $L_4$, the right side after substitution is convex in $x$.
For fixed $x$, it is the sum of a constant, $\frac{1}{L_4}$, and a
linear function of $L_4$, so it is convex in $L_4$ as well. The
endpoint bound in one variable, followed by the bound in the other, shows that
the expression is bounded by its largest corner value on each rectangle.

For the first rectangle, $\frac{A}{x^2(x-0.72)}$ equals $120$ at
$x=0.75$ and is less than $12$ at $x=0.93$. For the second rectangle,
the table for $\Psi$ supplies $31$ and $8$ as upper bounds at
$x=0.80$ and $x=0.93$. The table below records the resulting corner
estimates. The last two columns refer to the left and right endpoints
of the indicated interval for $L_4$.
\[
\begin{array}{c|c|c|cc}
 L_4 & x & \text{upper bound function at }x
     & \text{left} & \text{right}\\ \hline
 {[17,45]}  &0.75&120&<0.891&<0.989\\
 {[17,45]}  &0.93&12 &<0.997&<0.974\\
 {[45,150]} &0.80&31 &<0.879&<0.993\\
 {[45,150]} &0.93&8  &<0.967&<0.985
\end{array}
\]
For instance, the second row at $L_4=17$ has upper bound
$0.93+\frac{1}{17}+\frac{17\cdot12}{25000}<0.997$.
Every corner value is less than $0.998$, contrary to
\eqref{eq:wide-rectangle}. It follows that $L_4<17$.

We finally convert this bound on $L_4$ into a volume bound. For
$0<t<\mu_4$ with $t\le1$ there is a hollow translate of $tQ_4$.
Its longest chord in the direction defining $Q_3$ has length $tL_4$.
If $tL_4>1$, let $E$ be the set of projected points whose chords have
length at least one. By \cref{aux:chords}, $E$ is hollow and contains
a translate of $\alpha Q_3$, where
$\alpha=t(1-\frac{1}{tL_4})=t-\frac{1}{L_4}$. If $tL_4\ge4$,
then $4(1-\frac{1}{tL_4})^3\ge4(\frac{3}{4})^3>1$, and therefore
\eqref{aux:support-basic} yields $t^4|Q_4|\le tL_4|E|$. Moreover,
\eqref{aux:nested} bounds $L(E)$ by $\max\{tL_4,2tL_4-2\}=2(tL_4-1)$.
For $\alpha>0.6$, the first two bounds of \cref{aux:hollow-envelope}
applied to $E$ consequently imply
\[
 |Q_4|\le\frac{L_4}{t^3}
 \max\left\{\frac{729}{128},\frac{A\alpha}{\alpha-0.6},2A(tL_4-1)\right\}.
\]
If $L_4\le8$, then $8\mu_4>8(0.81)>6$. From \cref{aux:torus} with
$T=8$, $j=6$, we obtain $|Q_4|\le\frac{8^4}{6}<1700$.
If $8\le L_4<17$, take $t=0.81$. We have $tL_4>4$ and
$\alpha\ge0.81-\frac{1}{8}=0.685$. Since
$\frac{A(0.685)}{0.685-0.6}<52$ and $2A(0.81\cdot17-1)<52$,
the maximum in the last display is less than $52$. Here we used
that $\frac{\alpha}{\alpha-0.6}$ decreases with $\alpha$.
Thus $|Q_4|<\frac{17\cdot52}{(0.81)^3}<1700$ also in this range.
By \eqref{eq:wide-loss-direct},
$1-\mu_4<0.002+\frac{1700}{25000}=0.07$.

We can now use the larger scale $t=0.93$. If $L_4\le5$, then
$5\mu_4>4$, so \cref{aux:torus} with $T=5$ and $j=4$ yields
$|Q_4|\le\frac{5^4}{4}<200$. If $5\le L_4\le7$, then $tL_4>4$ and
$\alpha\ge0.93-\frac{1}{5}=0.73$. Since
$\frac{A(0.73)}{0.73-0.6}<22.5$ and $2A(0.93\cdot7-1)<22.5$,
we obtain $|Q_4|\le\frac{7(22.5)}{(0.93)^3}<200$.
For $7\le L_4\le17$, we have
$\alpha\ge0.93-\frac{1}{7}=\frac{551}{700}>\frac{16}{21}$,
so $\gamma\alpha>1$. By \eqref{aux:eq-hollow-envelope},
$|Q_4|\le\frac{L_4}{t^3}H_\gamma(\alpha)$. Since
$\frac{1}{L_4}=t-\alpha$, it remains to show
\[
 H_\gamma(\alpha)<200t^3(t-\alpha),\qquad
        \frac{551}{700}\le\alpha\le\frac{1481}{1700}.
\]
The left side is convex by \cref{aux:hollow-envelope}, and the right
side is affine. It suffices to check the two endpoints. For clarity,
we display strict upper bounds for the three terms of
$\frac{L_4}{t^3}H_\gamma(\alpha)$ at those endpoints, with $t=0.93$.
\[
\begin{array}{c|c|ccc}
 L_4&\alpha&
 \frac{L_4}{t^3}\frac{729}{128}&
 \frac{L_4}{t^3}\frac{A\alpha}{\alpha-0.6}&
 \frac{L_4}{t^3}\frac{A(\gamma\alpha)^3}{3(\gamma\alpha-1)}
 \\ \hline
 7 & \frac{551}{700} &50&75&196\\[2pt]
 17& \frac{1481}{1700} &121&138&149
\end{array}
\]
All six bounds are less than $200$. Thus $|Q_4|<200$ throughout the
last range as well. As observed at the start of the proof, this
contradicts $\mu_4\le0.99$ and proves the proposition.
\end{proof}

\section{The volume bound and its equality case}\label{sec:volume-proof}

The preceding three sections prove \cref{prop:projections} in all
its cases. We now apply \cref{lem:reduction} to the fiber over an
interior lattice point. The induction proves the volume bound and determines the
length of an edge at equality.

\begin{proof}[Proof of \cref{thm:general-rational}]
We argue by induction on $d$ for all the indicated values of $u$
simultaneously. The $k$ interior lattice points of an interval in
dimension one are consecutive integers. The first and last are $k-1$ apart,
and each is at distance at most one from its respective endpoint.
The interval therefore has length at most $1+(k-1)+1=k+1$. Here $A_1(u)=1$, and an interval
attaining the bound is itself an edge of length
$\frac{(k+1)f_0(u)}{u}=k+1$. This proves the initial case.

Suppose $d\ge2$. It suffices to consider $|C|\ge(k+1)A_d(u)$.
A lattice translation preserves volume, lattice-point counts and all
vertex denominators. Since $k\ge1$, we may therefore suppose
$0\in\intr C$.

We first consider $d=2$ and $w=\lw(C)\ge1+\frac{2}{u}$. Here
$u\ge2$. If $w\le2$, then \cref{aux:planar} and the monotonicity
of $\Phi$ imply
\[
 \frac{|C|}{k+1}\le\Phi(w)\le\Phi\left(1+\frac{2}{u}\right)
 =\frac{(u+2)^2}{4u}<\frac{(u+1)^2}{2u}=A_2(u).
\]
The strict comparison follows from
$2(u+1)^2-(u+2)^2=u^2-2>0$. If $w>2$, then \cref{aux:planar}
implies
\[
 \frac{|C|}{k+1}\le\frac{2k+\frac{5}{2}}{k+1}
 =2+\frac{1}{2(k+1)}\le\frac{9}{4}\le A_2(u).
\]
The penultimate inequality is an equality only at $k=1$.
For the last one,
$A_2(u)-\frac{9}{4}=\frac{(2u-1)(u-2)}{4u}$, with equality exactly at
$u=2$. Thus we have the volume bound in this planar case, and its
only possible equality parameters are the stated exception $(2,2,1)$.

For the remaining cases we choose a lattice projection $\pi$ of rank
$n$ as follows. If $\lw(C)<1+\frac{2}{u}$, take a primitive integral height
attaining the width, so $n=1$. Otherwise $d\ge3$, so there is a projection $\pi$ of rank $n\ge2$ with
$|\pi(C)|\le V_n(u)$ by \cref{prop:projections}.
In the latter case the image has width at least $\lw(C)\ge1+\frac{2}{u}$. In both cases the central fiber $C_0$ has dimension $r=d-n$,
and \cref{lem:reduction} bounds its vertex denominators by $U=f_n(u)$.
Since $0\in\intr C$, the relative interior of $C_0$ contains the
origin. By \eqref{eq:relative-fiber}, each of its lattice points lies
in $\intr C$. Thus $1\le I(C_0)\le k$.

We check that the theorem in dimension $r$ applies with the new
denominator bound $U$. If $u=n=1$, then $U=f_1(1)=2$ and $r=d-1\ge3$,
because the hypotheses with $u=1$ require $d\ge4$. If $u=1$ and
$n\ge2$, then $U\ge f_2(1)=6$. For $u\ge2$ we have
$U\ge f_1(u)=u(u+1)\ge6$. In all cases $1\le r<d$ and
$(r,U,I(C_0))\ne(2,2,1)$. Thus both the volume bound and the
assertion about an edge are available for $C_0$.

By the induction hypothesis and the definition of $A_r(U)$,
\[
 \Vol_r(C_0)=r!|C_0|
     \le \frac{(I(C_0)+1)f_{r-1}(U)^2}{U^{r+1}}.
\]
Since $r+n=d$ and $U=f_n(u)$, the recurrence implies
$f_{r-1}(U)=f_{d-1}(u)$. Substitution in \eqref{eq:one-lift}
cancels the factor $U^{r+1}$, and we obtain
\[
 \Vol_d(C)
 \le\frac{U^{r+1}}{u^{d+1}}\Vol_r(C_0)
 \le\frac{(I(C_0)+1)f_{d-1}(u)^2}{u^{d+1}}
 \le\frac{(k+1)f_{d-1}(u)^2}{u^{d+1}} .
\]
After division by $d!$ this is the required upper bound.

At equality every inequality in the last display is an equality.
In particular $I(C_0)=k$, and $C_0$ attains its own volume bound.
The induction hypothesis therefore shows that $C_0$ is an
$r$-simplex with an edge of lattice length
$\frac{(k+1)f_{r-1}(U)}{U}$. Since the pair $(d,u)=(2,1)$ is
outside the hypotheses, the equality assertion of
\cref{lem:reduction} applies. A face of $C$ is a translate of
$\frac{U}{u}C_0$, and this face together with $n$ further points
spans $C$. The face has $r+1$ vertices, so $C$ is the convex hull of
at most $r+1+n=d+1$ points. Its dimension $d$ forces these points to
be affinely independent. Hence $C$ is a simplex.

The corresponding edge of the face is an edge of $C$. A primitive
vector of the kernel lattice remains primitive in $N$. Indeed, if
it were an integral multiple $m z$ with $z\in N$ and $m>1$, then
$m\pi(z)=0$ in the quotient lattice, so $\pi(z)=0$, contradicting
primitivity in the kernel lattice. It follows that the lattice length
of the edge after dilation is
\[
 \frac{U}{u}\cdot\frac{(k+1)f_{r-1}(U)}{U}
 =\frac{(k+1)f_{d-1}(u)}{u} .
\]

The simplex following \cref{thm:general-rational} attains the bound
for every admissible $d,u,k$.
\end{proof}

\begin{proof}[Proof of \cref{thm:general-zpw}]
By \cref{thm:general-rational} with $u=1$, we have
$\Vol_d(P)\le(k+1)(s_d-1)^2$. Suppose equality holds. Then $P$
has an edge of lattice length $(k+1)(s_d-1)$.

The lattice diameter of $P$ is the maximum number of collinear lattice
points in $P$ minus one. Averkov, Kr\"umpelmann, and
Nill~\cite[Theorem~2.8 and the discussion of~(2.5)]{AKN} show that
the lattice diameter is at most $(k+1)(s_d-1)$,
with equality precisely when $P$ is affinely unimodularly equivalent
to $S_{d,k}$. The edge just obtained contains $(k+1)(s_d-1)+1$ lattice
points, including its endpoints, so the lattice diameter attains the bound. Conversely, $S_{d,k}$ has $k$ interior lattice points and
attains the volume bound.
\end{proof}



\vspace{5mm}
\noindent
\textbf{Acknowledgments.}
Feihu Liu was partially supported by the Postdoctoral Fellowship Program and China Postdoctoral Science Foundation (Grant No. BX2026002).

\noindent\textbf{Declaration of AI assistance.}
During the preparation of this manuscript, the authors used ChatGPT to assist in exploring possible approaches, checking technical details, and improving the exposition. All mathematical arguments, computations, proofs, and results were independently verified by the authors. The authors take full responsibility for the content of this manuscript.


\begin{thebibliography}{99}

\bibitem{AKN}
G.~Averkov, J.~Kr\"umpelmann, and B.~Nill,
\href{https://doi.org/10.1016/j.aim.2014.12.035}{\emph{Largest integral simplices with one interior integral point:
Solution of Hensley's conjecture and related results}},
Advances in Mathematics \textbf{274} (2015), 118--166.

\bibitem{AW}
G.~Averkov and C.~Wagner,
\href{https://doi.org/10.1007/s13366-011-0028-8}{\emph{Inequalities for the lattice width of lattice-free convex sets
in the plane}},
Beitr\"age zur Algebra und Geometrie \textbf{53} (2012), no.~1, 1--23.

\bibitem{BK}
G.~Balletti and A.~M.~Kasprzyk,
\emph{Three-dimensional lattice polytopes with two interior lattice points},
\href{https://arxiv.org/abs/1612.08918}{arXiv:1612.08918} (2016).

\bibitem{BKN}
G.~Balletti, A.~M.~Kasprzyk, and B.~Nill,
\href{https://doi.org/10.1017/fms.2022.93}{\emph{On the maximum dual volume of a canonical Fano polytope}},
Forum of Mathematics, Sigma \textbf{10} (2022), Article e109.

\bibitem{Bohnert}
M.~Bohnert,
\emph{Area bounds for planar convex bodies containing a fixed number
of interior integral points},
\href{https://arxiv.org/abs/2305.11485v2}{arXiv:2305.11485v2} (2025).

\bibitem{Gardner}
R.~J.~Gardner, \href{https://doi.org/10.1090/S0273-0979-02-00941-2}{\emph{The Brunn--Minkowski inequality}},
Bulletin of the American Mathematical Society (N.S.) \textbf{39}
(2002), no.~3, 355--405.

\bibitem{Hensley}
D.~Hensley, \href{https://msp.org/pjm/1983/105-1/pjm-v105-n1-p10-s.pdf}{\emph{Lattice vertex polytopes with interior lattice points}},
Pacific Journal of Mathematics \textbf{105} (1983), no.~1, 183--191.

\bibitem{Hurkens}
C.~A.~J.~Hurkens, \href{https://doi.org/10.1016/0024-3795(90)90010-A}{\emph{Blowing up convex sets in the plane}},
Linear Algebra and its Applications \textbf{134} (1990), 121--128.

\bibitem{Kasprzyk}
A.~M.~Kasprzyk, \href{https://doi.org/10.4153/CJM-2010-070-3}{\emph{Canonical toric Fano threefolds}},
Canadian Journal of Mathematics \textbf{62} (2010), no.~6, 1293--1309.

\bibitem{LZ}
J.~C.~Lagarias and G.~M.~Ziegler, \href{https://doi.org/10.4153/CJM-1991-058-4}{\emph{Bounds for lattice polytopes
containing a fixed number of interior points in a sublattice}},
Canadian Journal of Mathematics \textbf{43} (1991), no.~5, 1022--1035.

\bibitem{LTTZ}
F.~Liu, J.~Tang, S.~Tao, and Z.~Zhang,
\emph{The Zaks--Perles--Wills volume conjecture I:
three-dimensional lattice polytopes}, \href{https://arxiv.org/abs/2611.22337}{arXiv:2611.22337} (2026).

\bibitem{LTTZII}
F.~Liu, J.~Tang, S.~Tao, and Z.~Zhang,
\emph{The Zaks--Perles--Wills volume conjecture II: lattice simplices},
preprint, 2026.

\bibitem{Macbeath}
A.~M.~Macbeath,
\href{https://doi.org/10.1017/S0305004100028012}{\emph{On measure of sum sets II. The sum-theorem for the torus}},
Proceedings of the Cambridge Philosophical Society \textbf{49}
(1953), no.~1, 40--43.

\bibitem{Rockafellar}
R.~T.~Rockafellar, \href{https://doi.org/10.1515/9781400873173}{\emph{Convex Analysis}},
Princeton Mathematical Series, vol.~28, Princeton University Press,
Princeton, NJ, 1970.

\bibitem{ZPW}
J.~Zaks, M.~A.~Perles, and J.~M.~Wills,
\emph{On lattice polytopes having interior lattice points},
Elemente der Mathematik \textbf{37} (1982), no.~2, 44--46.

\end{thebibliography}
\end{document}
